\documentclass[11pt,letterpaper]{article}
\usepackage[margin=0.93in,headheight=15pt]{geometry}
\usepackage[T1]{fontenc}
\usepackage[utf8]{inputenc}
\usepackage{lmodern}
\usepackage{amsmath,amssymb,amsthm,mathtools}
\usepackage{microtype,booktabs,array,longtable,enumitem}
\usepackage{xcolor}
\definecolor{inkblue}{RGB}{27,62,94}
\usepackage{xurl}
\usepackage[colorlinks=true,linkcolor=inkblue,citecolor=inkblue,urlcolor=inkblue]{hyperref}
\usepackage{fancyhdr}
\pagestyle{fancy}
\fancyhf{}
\fancyhead[L]{\small Diophantine geometry in iPRA}
\fancyhead[R]{\small Explicit arithmetic proofs}
\fancyfoot[C]{\thepage}
\renewcommand{\headrulewidth}{0.35pt}
\setlength{\parskip}{0.35em}
\setlength{\parindent}{0pt}
\setlength{\emergencystretch}{2.5em}
\setlist{itemsep=0.25em,topsep=0.4em}
\newcommand{\N}{\mathbb N}
\newcommand{\Z}{\mathbb Z}
\newcommand{\Q}{\mathbb Q}
\newcommand{\R}{\mathbb R}
\newcommand{\PP}{\mathbb P}
\newcommand{\iPRA}{\ensuremath{\mathsf{iPRA}}}
\newcommand{\Ht}{\mathrm H}
\newcommand{\dotminus}{\mathbin{\dot{-}}}
\newtheorem{theorem}{Theorem}[section]
\newtheorem{lemma}[theorem]{Lemma}
\newtheorem{proposition}[theorem]{Proposition}
\newtheorem{corollary}[theorem]{Corollary}
\theoremstyle{definition}
\newtheorem{definition}[theorem]{Definition}
\newtheorem{example}[theorem]{Example}
\theoremstyle{remark}
\newtheorem{remark}[theorem]{Remark}
\numberwithin{equation}{section}
\pdfstringdefDisableCommands{\def\iPRA{iPRA}\def\Q{Q}\def\PP{P}\def\N{N}}
\hypersetup{pdftitle={Diophantine geometry in intuitionistic primitive recursive arithmetic},
 pdfauthor={Research report prepared for Vladimir Reshetnikov},
 pdfsubject={Heights, conics, explicit descent, arithmetic dynamics, and Lean formalization}}

\title{\textbf{Diophantine geometry in intuitionistic\\primitive recursive arithmetic}\\[0.65em]
 \large Heights, rational conics, explicit descent, and finite point algorithms}
\author{Research report prepared for Vladimir Reshetnikov}
\date{November 15, 2026}

\begin{document}
\maketitle
\thispagestyle{empty}

\begin{abstract}
We develop a finite-data theory of rational Diophantine geometry in
first-order intuitionistic primitive recursive arithmetic. Rational projective
points are canonical primitive integer vectors. We prove exact bounded-height
enumeration, rational-root bounds, and two-sided height inequalities for
nondegenerate maps of the projective line using explicit Sylvester-adjugate
identities. These give a primitive recursive algorithm listing every rational
preperiodic point of a map of degree at least two. For a nonsingular conic
with a supplied rational point, we construct a total parametrization,
including the tangent parameter, and prove quadratic height comparison and
linear point-count bounds. A fully specified finite descent proves Fermat's
square-area theorem and determines all rational points of
$y^2=x^3-x$. Further results give congruence certificates, square and cube
denominators on Weierstrass equations, and a constructive finite-generation
interface with an explicit reducing selector.

All the stated arithmetic results have human-readable proofs with bounded
searches or specified primitive recursions and quantifier-free induction
invariants. Thirteen arithmetic-core theorems are checked in Lean 4.32.0
using only \texttt{Std}. The package includes exact executable companions,
validation evidence, and a detailed dependency plan for constructing the
remaining deeply embedded iPRA derivations in ProveIt. No complete
object-theory Lean formalization of the geometry is claimed.
\end{abstract}

\vfill
\begin{center}
\small\textbf{Verification status:} full written proofs for the stated scope;
checked Lean arithmetic core; explicit plan for object-theory completion.
\end{center}
\clearpage
\begingroup
\small
\setlength{\parskip}{0pt}
\tableofcontents
\endgroup
\clearpage

\section{What the development proves, and the arithmetic it uses}
\label{sec:foundations}

Diophantine geometry studies polynomial equations through their rational or
integral points, coordinate changes, heights, and descent. A first-order
arithmetic can express a substantial part of this subject if these objects are
represented by finite data. This paper develops that part over $\Q$. The central
results are an exact bounded-height point enumerator; explicit height distortion
under a nondegenerate map of $\PP^1$; a complete algorithm for its rational
preperiodic points; parametrization and quantitative point counts for a conic
with a supplied point; and the complete calculation
\[
   \{(x,y)\in\Q^2:y^2=x^3-x\}
       =\{(-1,0),(0,0),(1,0)\}.
\]
The last result follows from a fully specified finite descent. We also develop
congruence certificates, rational-root bounds, denominators on integral
Weierstrass equations, and a constructive finite-generation descent interface.

These are a selected theory, not a claim to reproduce the entire subject.
In particular the paper does not prove the general Mordell--Weil, Roth, Siegel,
or Faltings theorems. It does not assert their unprovability in $\iPRA$ either.
Each positive theorem below has its own finite arithmetic proof. Classical
background and terminology may be compared with Milne \cite{milne}, Sutherland
\cite{sutherland2013,sutherland2025}, and Silverman \cite{silverman}.

\subsection{The precise base theory}

We use first-order intuitionistic primitive recursive arithmetic, with one
sort $\N$ and equality. There is a symbol for each fixed primitive recursive
definition: zero, successor, projections, composition, and primitive recursion.
The projection symbols satisfy $\pi_i^k(x_0,\ldots,x_{k-1})=x_i$;
the composition symbol for $f$ and $g_1,\ldots,g_r$ satisfies
$c(\vec x)=f(g_1(\vec x),\ldots,g_r(\vec x))$.
For $h(0,\vec x)=f(\vec x)$ and
$h(n+1,\vec x)=g(n,h(n,\vec x),\vec x)$, these two equations are axioms.
The nonlogical axioms also include successor injectivity and $Sx\ne0$.
Induction is the rule
\[
 \frac{\Gamma\vdash A(0,\vec u)\qquad
       \Gamma,A(n,\vec u)\vdash A(n+1,\vec u)}
      {\Gamma\vdash A(t,\vec u)},
 \qquad A\text{ quantifier-free},
\]
with the usual freshness condition on $n$. The logical rules are intuitionistic
natural deduction: introduction and elimination for $\wedge,\vee,\to,\forall,
\exists$, equality substitution, and elimination of $\bot$. There is no
excluded-middle axiom for arbitrary formulas and no induction on arbitrary
formulas. Quantifier-free formulas are decidable because their atomic
equalities have primitive recursive characteristic functions.

Integers, rationals, vectors and polynomial lists are definitional encodings
in this single numerical sort. For example encode $z\ge0$ by $2z$ and
$z<0$ by $-2z-1$. A rational is a pair $(a,b)$ with $a\in\Z$, $b\ge1$,
$\gcd(|a|,b)=1$; addition and multiplication reduce by gcd. Finite sequences
can be encoded by a length and an iterated pairing, with primitive recursive
access functions. Invalid codes have fixed default outputs. Every geometric
assertion carries the corresponding validity predicate.

\subsection{Bounded quantifiers and explicit witnesses}

A bounded search is a recursion through $0,\ldots,B$. It retains the first
successful index and a flag, or reports failure after exactly $B+1$ tests.
A bounded conjunction is computed by Boolean multiplication. Thus a formula
whose quantifiers are all explicitly bounded has a primitive recursive
characteristic function. Its induction is an instance of quantifier-free
induction applied to that function.

Here is the usual elementary justification, including the point relevant to
weak arithmetic. A search state stores either no successful index, or an
index $j\le i$ for which the tested equality is true. The invariant that the
stored index is valid is quantifier-free. For each additional parameter $k$,
the invariant that a successful test at $k\le i$ has set the flag is also
quantifier-free. Induction on $i$, followed by universal introduction on $k$,
proves success and failure correctness. Bounded conjunction is handled by
searching for the first failed test. Nested bounded quantifiers are compiled
one layer at a time. This explains the use of finite maxima, finite
pigeonhole, and bounded lists later without strengthening induction.

There is an essential distinction between a fixed primitive recursive
definition and a variable code for an arbitrary primitive recursive program.
We use fixed algorithms with numerical inputs, including variable dimensions,
degrees and coefficient lists. We do \emph{not} assume a primitive recursive
universal evaluator for all primitive recursive programs.

\subsection{Elementary arithmetic infrastructure}

\begin{lemma}[Finite arithmetic operations]\label{lem:arithmetic}
The following operations and their stated elementary properties are available
in $\iPRA$: integer arithmetic and order; division with remainder by a positive
integer; gcd and Bezout coefficients; bounded divisibility tests; integer
square root; finite products; and prime factorization of a positive integer
by trial division. Coprime factors of a square are squares. If $a,b$ are
coprime, so are $a^r,b^s$ for positive $r,s$.
\end{lemma}
\begin{proof}
Division of $n$ by $d>0$ searches $q\le n$ for the largest $dq\le n$ and sets
$r=n-dq$. Maximality gives $r<d$. The equations $n=dq+r$ and $0\le r<d$
give uniqueness by subtraction. The square-root algorithm searches
$s\le n$ for the largest $s^2\le n$; it gives
$s^2\le n<(s+1)^2$.

Euclid's remainder algorithm can be run for $a+b+1$ stages, keeping the state
unchanged after a zero remainder. Every active stage strictly reduces the
positive second entry, so it stops in this bound. Maintaining two integer
linear-combination rows proves that the terminal positive entry $g$ is an
integer combination of $a,b$. The identity
$\gcd(a,b)=\gcd(b,a-qb)$ follows directly from the two divisibility
implications. It proves that $g$ divides both inputs and every common divisor
divides $g$. Iterating this construction gives Bezout coefficients for a
nonzero finite vector. The case of zero inputs is handled separately.

For factorization, visit the integers $p=2,\ldots,n$, repeatedly divide the
current residual by $p$ for at most $n$ stages, and record the exponent.
The residual-product invariant is an equation. After processing $p$, no
prime at most $p$ divides the residual: a nontrivial repeated division
strictly decreases it, so $n$ stages suffice to remove that divisor.
This invariant is a bounded test. An active divisor is prime:
if it had a smaller prime divisor, that divisor would already have been
removed. Existence of a prime divisor of a number $>1$ is obtained by its
least divisor in the bounded interval $[2,n]$; a factorization of that divisor
would contradict minimality. Euclid's prime-divisor lemma follows from
Bezout: if $p\nmid a$ then $up+va=1$, and multiplication by $b$ shows
$p\mid ab\Rightarrow p\mid b$. This also proves uniqueness of the exponent
of each prime in a product. If coprime $a,b>0$ have square product, each prime
occurs in just one of them, with even exponent; multiplying the corresponding
half powers constructs their square roots. The same argument proves
coprimality of powers. All divisions, trial indices and exponents are bounded
by the input integer or its explicitly given product. All loop invariants
just used are equations or bounded divisibility assertions, hence Boolean
quantifier-free invariants as above.
\end{proof}

\begin{lemma}[Primitive representatives]\label{lem:normalize}
Every nonzero integral vector $v=(v_0,\ldots,v_n)$ has a primitive
representative $\operatorname{norm}(v)$, obtained by dividing by
$g=\gcd(|v_0|,\ldots,|v_n|)$ and choosing the sign so the first nonzero
coordinate is positive. Two such representatives of the same rational
projective point are equal. Normalization is primitive recursive and never
increases the maximum absolute coordinate.
\end{lemma}
\begin{proof}
The construction and the bound follow from $g\ge1$. Suppose primitive
vectors $v,w$ satisfy $v_iw_j=v_jw_i$ for all $i,j$. Choose an index $j$
with $v_j\ne0$, and Bezout coefficients $c_i$ with $\sum_i c_iv_i=1$.
Set $a=\sum_i c_iw_i$. The proportionality equations give $v_ja=w_j$ and
then $v_jw_i=v_ja v_i$, so cancellation gives $w_i=av_i$ for every $i$.
As $w$ is primitive, $|a|=1$. The sign convention forces $a=1$.
Clearing the finitely many denominators of any nonzero rational vector
reduces to the integral case. Choosing $j$ and summing the Bezout equations
are bounded operations.
\end{proof}

We therefore use canonical primitive vectors as the actual codes for
$\PP^n(\Q)$. Quotient types and choices of representatives are unnecessary
in the object theory. A homogeneous polynomial is a finite list of integer
coefficients and exponent vectors of one declared total degree. The zero
polynomial is allowed with any declared degree; subsequent degree bounds
refer to that declared degree. Its value is computed
by bounded products and sums. Homogeneity shows that its zero locus is
unchanged by scaling and normalization.

\subsection{Relation to the inspected ProveIt foundation}

The inspected repository revision is
\texttt{b5a3ad037723c44d69a3d5b28a950c554327a989}.
The files \texttt{PR.lean}, \texttt{BiInterpretation.lean}, and
\texttt{NativeBiInterpretation.lean} were read at that revision, together
with the ConciseIPRA README. The first defines the first-order theory just
described, with function-symbol type \texttt{PR}, term type \texttt{PTm},
and axiom family \texttt{PRAAx}. The other two expose the translation to the
scalar and native presentations. In particular
\texttt{pra\_iff\_native} transfers a derivation of an iPRA formula to a
derivation of its native reading \cite{proveit}.

This is the appropriate foundation to reuse. The new geometry has not yet
been turned into values of its derivation type. The standalone Lean core in
this package verifies selected arithmetic lemmas; Section~\ref{sec:lean}
specifies the additional work needed for deeply embedded derivations. The
repository's archived-draft index was inspected, but no theorem from an
unread archive is used as a proof dependency.

\section{Heights, finite point lists, and rational roots}
\label{sec:heights}

\subsection{An exact arithmetic form of Northcott finiteness over \texorpdfstring{$\Q$}{Q}}

For a canonical primitive vector $x$, define the multiplicative height
\[
                 \Ht([x_0:\cdots:x_n])=\max_i |x_i|.
\]
It is a positive integer. For $r=a/b$ in lowest terms with $b>0$,
$\Ht(r)=\max(|a|,b)$, by identifying $r$ with $[a:b]$.
We use no logarithms or real-valued heights.

\begin{theorem}[Bounded-height enumeration]\label{thm:northcott}
There is a primitive recursive algorithm $\mathcal L(n,B,F_1,\ldots,F_s)$
whose output is a duplicate-free list of exactly the points
\[
 \{P\in\PP^n(\Q):\Ht(P)\le B,\ F_1(P)=\cdots=F_s(P)=0\},
\]
for homogeneous integral input polynomials. Its length is at most
$(2B+1)^{n+1}$. Its correctness is provable in $\iPRA$.
\end{theorem}
\begin{proof}
Visit each vector in $[-B,B]^{n+1}$ in a fixed lexicographic order. Retain
it exactly when its gcd is one, its first nonzero entry is positive, and all
polynomial values vanish. There are precisely $(2B+1)^{n+1}$ visits. Every
retained vector is a point of the required height on the locus. Conversely
the canonical vector of any required point belongs to this box and passes
every test. Uniqueness is Lemma~\ref{lem:normalize}.

For $B=0$ the only visited vector is zero, which is rejected. For $B\ge1$,
an explicit numerical implementation uses a visit counter with base
$(2B+1)$ digits of length $n+1$; subtract $B$ from each digit. The output
list and its length are maintained by a bounded recursion. Polynomial
evaluation on the box satisfies
\[
    |F(x)|\le \|F\|_1\max(1,B)^{\deg F},\qquad
    \|F\|_1=\sum_{\alpha}|c_\alpha|,
\]
because each monomial is bounded by the corresponding power. This estimate
also bounds coefficient-weighted monomials and their partial sums. A raw
monomial computation is bounded by $\max(1,B)^{\deg F}$; list codes are
bounded separately by iterating the fixed pairing function for the stated
number of visits. With the input point as an extra parameter,
the invariant ``if its visit has occurred and it passes the tests, it is in
the retained list'' is bounded. This proves completeness using the Boolean
induction of Section~\ref{sec:foundations}.
\end{proof}

The theorem is the rational, fixed-field form of the height finiteness
principle. It supplies a complete list below a \emph{given} height, not an
upper bound for all points of every variety. Those are different assertions.
When $B=0$ the list is empty, as it must be.

\begin{corollary}[Effective use of a global height bound]\label{cor:heightbound}
For a homogeneous projective polynomial system $V$, suppose an integer $B$ and a proof of
$P\in V(\Q)\Rightarrow\Ht(P)\le B$ are supplied. Then
$\mathcal L(n,B,V)$ is the complete rational-point list. In particular,
its emptiness decides whether $V$ has a rational point.
\end{corollary}
\begin{proof}
Soundness is the previous theorem. The supplied implication reduces every
point to its completeness clause. No search for $B$ is part of the argument.
\end{proof}

\subsection{Heights under displayed coordinate maps}

\begin{proposition}[Upper height bounds]\label{prop:upperheight}
Let $F=(F_0,\ldots,F_m)$ be homogeneous forms of degree $d$ and suppose
the evaluated vector is nonzero at $P$. Put
$A=\max_i\|F_i\|_1$. Then
\[
                 \Ht(F(P))\le A\Ht(P)^d.
\]
If an integral square matrix $M$ has nonzero determinant and row norm
$\|M\|_\infty=\max_i\sum_j|M_{ij}|$, then
\[
 \Ht(MP)\le\|M\|_\infty\Ht(P),\qquad
 \Ht(P)\le\|\operatorname{adj}M\|_\infty\Ht(MP).
\]
\end{proposition}
\begin{proof}
The monomial estimate above bounds the unnormalized image coordinates.
Normalization only decreases them. For the inverse bound, apply the first
inequality to the integral adjugate and use
$\operatorname{adj}(M)M=\det(M)I$. Multiplication by the nonzero scalar
$\det(M)$ fixes the projective point. The adjugate entries are finite
determinant sums, so the constants are explicit primitive recursive data.
\end{proof}

For completeness, the adjugate identity has an elementary finite proof.
The diagonal entry is a cofactor expansion of the determinant. An off-diagonal
entry is a determinant with two equal rows, which is zero: pair each
permutation summand with the summand obtained by exchanging the two row
assignments. Their products agree and their signs are opposite. This is a
bounded involution on a list of at most $k!$ permutations for a $k\times k$
matrix, so its proof is available in $\iPRA$ by finite sums. This is also the
determinant fact used in the next section.

\subsection{A complete zero-dimensional example: rational roots}

\begin{theorem}[Rational-root bounds and a complete list]\label{thm:roots}
Let $f(T)=a_0+\cdots+a_dT^d$ be nonzero, with $a_d\ne0$, and let
$k$ be the least index with $a_k\ne0$. Every nonzero rational root $p/q$
in lowest terms, $q>0$, satisfies
\[
                  p\mid a_k,\qquad q\mid a_d.
\]
It therefore has height at most $\max(|a_k|,|a_d|)$. All rational roots
are obtained by the finite divisor search, with $0$ included exactly when
$a_0=0$. This search and its correctness are primitive recursive and
provable in $\iPRA$.
\end{theorem}
\begin{proof}
Remove the factor $T^k$ and write $g(T)=b_0+\cdots+b_eT^e$ with
$b_0=a_k\ne0$ and $b_e=a_d\ne0$. For $e=0$ there is no nonzero root.
For $e>0$, clear denominators in $g(p/q)=0$:
\[
       b_0q^e+b_1pq^{e-1}+\cdots+b_ep^e=0.
\]
All but the first summand are divisible by $p$, so $p\mid b_0q^e$.
Coprimality of $p,q^e$ and Bezout give $p\mid b_0$. Likewise reducing
the equality modulo $q$ yields $q\mid b_ep^e$ and hence $q\mid b_e$.
The divisibility bounds follow because both coefficients are nonzero.
Enumerate $-|b_0|\le p\le|b_0|$, $p\ne0$, and
$1\le q\le|b_e|$, retain coprime pairs with these divisibilities, and
evaluate $f(p/q)$ exactly. Soundness is the test and completeness is the
argument just given. Finally $f(0)=a_0$ decides the zero root separately.
\end{proof}

For example $T^3-5T^2+6T$ has the complete list $0,2,3$. The zero
polynomial is deliberately excluded: every rational would be a root. A
nonzero constant has the empty list. These cases matter for a total
certificate checker.

\section{Explicit height functoriality and complete arithmetic dynamics}
\label{sec:dynamics}

A lower height bound must account for common factors removed from the image
coordinates. Bounding the displayed coordinates alone would not suffice.
The following determinant construction controls those factors and supplies
an effective finiteness theorem. Its classical setting is the height theory
of arithmetic dynamics \cite{silverman,silverman2011}; the proof here is
entirely in integer arithmetic.

\subsection{Two finite polynomial certificates}

Let
\[
 F(X,Y)=\sum_{i=0}^d f_iX^{d-i}Y^i,\qquad
 G(X,Y)=\sum_{i=0}^d g_iX^{d-i}Y^i,\qquad d\ge1.
\]
For a nonnegative integer $r$, a homogeneous binary form of degree $r$ is
represented by its $r+1$ coefficient list. Multiplication is finite
convolution. Form the $2d\times2d$ integer matrix $S$ of the map
\[
   (A,B)\longmapsto AF+BG,
\]
where $A,B$ have degree $d-1$ and the output has degree $2d-1$.
Use the ordered monomial basis $X^r,X^{r-1}Y,\ldots,Y^r$ in every space.
Let $R=\det S$. The nondegeneracy hypothesis for this section is the
decidable arithmetic condition $R\ne0$. This is the homogeneous resultant
condition, defined here by its matrix rather than by an algebraic-closure
criterion.

\begin{lemma}[Computed Bezout certificates]\label{lem:resultant}
The first and last columns of $\operatorname{adj}S$ provide forms
$A_X,B_X,A_Y,B_Y$ of degree $d-1$ with
\begin{align}
 A_XF+B_XG&=R X^{2d-1},\label{eq:bezout-X}\\
 A_YF+B_YG&=R Y^{2d-1}.\label{eq:bezout-Y}
\end{align}
These identities are provable in $\iPRA$ and are checked coefficientwise.
If
\[
 M=\max(1,|f_0|,\ldots,|f_d|,|g_0|,\ldots,|g_d|),
 \quad K=(2d-1)!M^{2d-1},
\]
every coefficient of the four certificate forms has absolute value at most
$K$.
\end{lemma}
\begin{proof}
The adjugate identity gives $S\operatorname{adj}S=RI$; its first and last
columns are precisely the two coefficient identities. Each cofactor is a
sum of $(2d-1)!$ signed products of $2d-1$ entries of $S$. Every entry of
$S$ is a coefficient of $F$ or $G$, or zero. Bounding each product by
$M^{2d-1}$ proves the claim. Determinants, permutations, signs and coefficient
convolution are finite sums and products with the displayed bounds; the
finite proof of the adjugate identity was given in
Section~\ref{sec:heights}.
\end{proof}

Set
\[
 L=\max\bigl(\|A_X\|_1+\|B_X\|_1,
             \|A_Y\|_1+\|B_Y\|_1\bigr),\qquad U=(d+1)M.
\]
Since $R\ne0$, $L\ge1$. The coarse bound
\begin{equation}\label{eq:coarse-L}
                L\le 2d(2d-1)!M^{2d-1}
\end{equation}
is useful if one does not wish to keep the smaller computed constant.

\begin{theorem}[Two-sided height functoriality on $\PP^1$]
\label{thm:height-map}
If $R\ne0$, the rule $f([x:y])=[F(x,y):G(x,y)]$ is a total map on
$\PP^1(\Q)$. For every rational projective point $P$,
\begin{equation}\label{eq:height-map}
              \Ht(P)^d\le L\Ht(f(P)),\qquad
              \Ht(f(P))\le U\Ht(P)^d.
\end{equation}
The map, constants, and proofs are uniform primitive recursive constructions
from the finite coefficient lists.
\end{theorem}
\begin{proof}
Take a primitive integral pair $(x,y)$ and put $H=\max(|x|,|y|)\ge1$.
If $F(x,y)=G(x,y)=0$, the two identities imply
$R x^{2d-1}=R y^{2d-1}=0$. Since $R\ne0$ and $2d-1>0$, this would give
$x=y=0$. Thus the image is nonzero. Homogeneity establishes that its
projective class is independent of the chosen representative.

Put $T=\max(|F(x,y)|,|G(x,y)|)$ and
$g=\gcd(|F(x,y)|,|G(x,y)|)\ge1$. The identities show that $g$ divides
both $R x^{2d-1}$ and $R y^{2d-1}$. The two powers are coprime; multiplying
their Bezout identity by $R$ shows $g\mid R$. Consequently $g\le|R|$.
This is the required control of projective cancellation.

Choose a coordinate of absolute value $H$ and use the corresponding
certificate. Each coefficient form has degree $d-1$, so
\[
 |R|H^{2d-1}\le L H^{d-1}T.
\]
Cancel the positive integer $H^{d-1}$ to obtain
$|R|H^d\le LT$. Since $T=g\Ht(f(P))$ and $g\le|R|$, we have
\[
 |R|H^d\le Lg\Ht(f(P))\le L|R|\Ht(f(P)).
\]
Cancellation of $|R|>0$ proves the lower bound. For the upper bound,
$T\le(d+1)MH^d=UH^d$, and division by $g$ cannot increase height.
All case distinctions here concern integer comparisons or finite lists.
\end{proof}

The proof also works with any supplied forms satisfying
\eqref{eq:bezout-X}--\eqref{eq:bezout-Y} and any nonzero $R$; it need not
recompute the adjugate. The same two finite identities prove nonvanishing,
bound the gcd, and give the lower height estimate. This is a useful
certificate interface for formalization.

\subsection{A global finiteness theorem with an actual decision procedure}

Define the iterate by $f^0(P)=P$ and $f^{n+1}(P)=f(f^n(P))$. A point is
\emph{preperiodic} if
\[
             \exists i\,\exists j\ (i<j\ \wedge\ f^i(P)=f^j(P)).
\]
This is initially an unbounded existential statement. The next theorem
reduces it to a bounded computation.

\begin{theorem}[Primitive recursive Northcott theorem for rational dynamics]
\label{thm:preperiodic}
Suppose $d\ge2$ and $R\ne0$. Every rational preperiodic point of $f$ has
height at most $L$. There is a primitive recursive algorithm that returns
\emph{all} rational preperiodic points, together with a repetition witness
for each, and whose correctness is provable in $\iPRA$.
In particular there are at most $(2L+1)^2$ such points.
\end{theorem}
\begin{proof}
If $H=\Ht(P)>L$, then
$H^d\ge H^2>LH$. By \eqref{eq:height-map}, $L\Ht(f(P))\ge H^d>LH$,
so $\Ht(f(P))>H$. Induction on a numerical iteration counter therefore
shows that all subsequent heights strictly increase. Given $i<j$, a
finite induction from $i$ through $j$ proves
$\Ht(f^i(P))<\Ht(f^j(P))$. Such a point cannot have a repetition witness.
Thus any preperiodic point has height at most $L$.

Every iterate of a preperiodic point is preperiodic: apply $f^t$ to a
witness equality and use $f^i(f^t(P))=f^t(f^i(P))$, an identity proved by
finite recursion. It follows that its entire forward orbit remains in
the explicit finite list
\[
             S=\mathcal L(1,L),\qquad N=|S|\le(2L+1)^2.
\]

The algorithm first constructs $S$. For each $P\in S$, start an empty
visited list and follow its orbit for at most $N+1$ inspected states. If
the current height exceeds $L$, reject $P$. If the current state equals
a previously visited state, retain $P$ and the two visit indices. Otherwise
append it and continue. The transition can be made a map on the bounded
set $S$ plus one escape flag: after escape, keep the flag unchanged.

No candidate can remain unclassified. If $N+1$ states were inspected
without escape or repetition, they would be $N+1$ different members of an
$N$-element list. Here is a finite arithmetic proof of the contradiction.
Find each state's index in $S$ by bounded search. For every $r<N$, let
$c_r$ count the occurrences of that index among the $N+1$ visits. Absence
of repetitions gives $c_r\le1$, hence $\sum_{r<N}c_r\le N$. On the
other hand each visit has exactly one index, so summing the same incidence
table by visits gives $\sum_{r<N}c_r=N+1$. Both equalities are bounded
double-sum identities.

Soundness of retention is the returned equality with two distinct indices.
If a candidate is preperiodic, all its iterates lie in $S$, so rejection
by escape cannot occur; the preceding pigeonhole argument then retains it.
The earlier height bound proves that no preperiodic point was omitted from
the candidate list.

The certificate construction, construction of $S$, and the loop counts are
bounded by primitive recursive functions of the coefficient lists. Before
escape, raw image coordinates are at most $UL^d$. Keeping the escape flag
prevents uncontrolled growth of the stored states. Thus the whole procedure
is a fixed primitive recursive algorithm, not a merely terminating search
whose running time has been left unspecified. Its invariants are Boolean
tests, list membership, and bounded counting equations.
\end{proof}

One may use the coarser right side of \eqref{eq:coarse-L} everywhere in
place of $L$ and still obtain a primitive recursive bound depending only
on $d,M$. The bound is usually far too large for practical computation.
Primitive recursive totality and computational efficiency are separate issues.

\subsection{Two complete examples}

For $f(t)=t^2$, use $F=X^2$, $G=Y^2$ and the identities
$XF=X^3$, $YG=Y^3$. Here $R=1$ and $L=1$.
The rational projective points of height at most one are
$\infty,0,1,-1$. Their images are respectively $\infty,0,1,1$.
All four are preperiodic, and Theorem~\ref{thm:preperiodic} proves that
there are no others.

For $f(t)=t^2-1$, use $F=X^2-Y^2$, $G=Y^2$ and
\[
                 XF+XG=X^3,\qquad YG=Y^3.
\]
Thus $R=1$ and $L=2$. The complete height-two list is
\[
             \infty,\ 0,\ \pm1,\ \pm2,\ \pm\tfrac12.
\]
The four points $\infty,0,1,-1$ are retained: $\infty$ is fixed,
$0\leftrightarrow-1$ is a cycle, and $1\mapsto0$.
Both $2$ and $-2$ map to $3$, whose height exceeds two. Both
$1/2$ and $-1/2$ map to $-3/4$, whose height exceeds two.
They cannot be preperiodic by the proven escape bound. Consequently
\[
          \operatorname{PrePer}(t^2-1,\Q)=\{\infty,0,1,-1\}.
\]
These calculations are reproduced by the exact Python companion. Their
completeness comes from the height proof, not from a numerical search cutoff.

\section{Conics: complete parametrization and explicit height growth}
\label{sec:conics}

The first substantial geometric application is a complete description of
the rational points on a nonsingular conic \emph{with a supplied rational
point}.  The geometric idea is projection from that point: a line through
it has one other intersection with the conic, with the tangent line
returning the original point.  This is the usual presentation of rational
points on a conic; see, for example, \cite{sutherland2025}.  We give all the
formulas, handle the tangent case explicitly, and derive quantitative
height estimates by integer identities.  Consequently neither a theorem
about algebraically closed fields nor an intersection-multiplicity
formalism is needed.

\subsection{A normal form that preserves explicit information}

Let
\[
 q(X_0,X_1,X_2)=\sum_{0\leq i\leq j\leq2}a_{ij}X_iX_j,
 \qquad a_{ij}\in\Z,
\]
and let
\[
 S_q=\begin{pmatrix}
 2a_{00}&a_{01}&a_{02}\\
 a_{01}&2a_{11}&a_{12}\\
 a_{02}&a_{12}&2a_{22}
 \end{pmatrix}.
\]
We use the decidable nonsingularity condition $\det S_q\ne0$.
It implies the usual assertion that the projective conic has no singular
point: its gradient at a vector $x$ is $S_qx$, and multiplication by the
adjugate matrix shows that $S_qx=0$ implies $x=0$.  All matrices here have
the fixed size three, so this argument is an identity of integer
polynomials.  Put
\[
 C=\sum_{i\leq j}|a_{ij}|.
\]
Nonsingularity implies $C\geq1$.

Suppose that a nonzero primitive integer vector
$p=(p_0,p_1,p_2)$ with $q(p)=0$ is given, and write
$P=\max_i|p_i|\geq1$.  A permutation of coordinates makes $p_0\ne0$.
This permutation is found by checking the three coordinates; it preserves
$C$ and $P$.  Define
\[
 M=\begin{pmatrix}p_0&0&0\\p_1&1&0\\p_2&0&1\end{pmatrix}.
\]
Since $\det M=p_0\ne0$, it induces a bijection of rational projective
planes.  Its inverse can be evaluated without rational division:
\begin{equation}\label{eq:conic-adjugate}
 \operatorname{adj}(M)x
 =(x_0,\ p_0x_1-p_1x_0,\ p_0x_2-p_2x_0).
\end{equation}
Indeed $M\operatorname{adj}(M)=p_0I$.

The coefficient of $U^2$ in $q(M(U,V,W))$ is $q(p)=0$.
Write the remaining coefficients as
\begin{equation}\label{eq:conic-normal}
 q(M(U,V,W))=U L(V,W)+R(V,W),
 \qquad
 L(V,W)=\ell V+mW,\quad R(V,W)=aV^2+bVW+cW^2.
\end{equation}
Explicitly,
\[
\begin{split}
 \ell&=a_{01}p_0+2a_{11}p_1+a_{12}p_2,\\
 m&=a_{02}p_0+a_{12}p_1+2a_{22}p_2,\\
 a&=a_{11},\qquad b=a_{12},\qquad c=a_{22}.
\end{split}
\]
The polar matrix of the transformed form is $M^{\mathsf T}S_qM$.
Expansion of its determinant gives
\begin{equation}\label{eq:conic-r}
 \det(M^{\mathsf T}S_qM)
 =\det\begin{pmatrix}0&\ell&m\\\ell&2a&b\\m&b&2c\end{pmatrix}
 =-2r,
 \qquad r=am^2-b\ell m+c\ell^2=R(-m,\ell).
\end{equation}
Since this determinant also equals $p_0^2\det S_q$, we have $r\ne0$.
In particular $(\ell,m)\ne(0,0)$.

The following two identities will serve both to exclude undefined cases
and to control the cancellation in projective coordinates:
\begin{align}
 rs^2&=m^2R(s,t)
   -L(s,t)\bigl((bm-c\ell)s+cmt\bigr),\label{eq:conic-identity-s}\\
 rt^2&=\ell^2R(s,t)
   -L(s,t)\bigl(a\ell s+(b\ell-am)t\bigr).
   \label{eq:conic-identity-t}
\end{align}
For example, expansion of the right side of
\eqref{eq:conic-identity-s} cancels its $st$ and $t^2$ terms and leaves
$(am^2-b\ell m+c\ell^2)s^2$.  Expanding the other identity gives the
same verification.  If $L(s,t)=R(s,t)=0$, these identities and $r\ne0$
imply $s=t=0$.  Thus $L$ and $R$ have no common projective zero over
$\Q$; this conclusion has been obtained without a resultant theorem.

\subsection{The parametrization, including its tangent value}

Write $C_q(\Q)$ for the rational projective points satisfying $q=0$.
For $(s,t)\ne(0,0)$ define
\begin{equation}\label{eq:conic-param}
 F(s,t)=\bigl(-R(s,t),\ sL(s,t),\ tL(s,t)\bigr),
 \qquad
 \Phi([s:t])=[M F(s,t)].
\end{equation}
All projective vectors are normalized by the primitive normalization
procedure when a canonical point code is required.

\begin{theorem}[Effective parametrization of a pointed conic]
\label{thm:conic-param}
Under the hypotheses above, $\Phi$ is a bijection
\[
 \Phi:\PP^1(\Q)\longrightarrow C_q(\Q).
\]
It and its inverse are primitive recursive on canonical point codes,
uniformly in the six coefficients and the supplied point $p$.
The parameter $[-m:\ell]$ maps to $[p]$.
\end{theorem}

\begin{proof}
First, $F(s,t)\ne0$.  If $L(s,t)\ne0$, at least one of its last two
coordinates is nonzero.  If $L(s,t)=0$, then $R(s,t)\ne0$ by
\eqref{eq:conic-identity-s}--\eqref{eq:conic-identity-t}.
Since $M$ is invertible, $MF(s,t)\ne0$ as well.  All coordinates of $F$
are homogeneous of degree two.  Rescaling $(s,t)$ by a nonzero rational
number rescales $F$ by its square; hence the formula is well defined
on projective points.

Substitution in \eqref{eq:conic-normal}, using homogeneity of $L$ and
$R$, gives
\[
\begin{aligned}
 q(MF(s,t))&=(-R(s,t))L(sL(s,t),tL(s,t))\\
           &\qquad+R(sL(s,t),tL(s,t))\\
           &=-R(s,t)L(s,t)^2+R(s,t)L(s,t)^2=0.
\end{aligned}
\]
Thus the image lies on the conic.

To construct the inverse, first transform a conic point $[x]$ by
\eqref{eq:conic-adjugate}, obtaining $[u:v:w]$ satisfying
\[
 uL(v,w)+R(v,w)=0.
\]
If $(v,w)\ne(0,0)$, then $L(v,w)\ne0$: otherwise this equation would
also give $R(v,w)=0$, contrary to the two identities above.  Therefore
\[
 F(v,w)=(-R(v,w),vL(v,w),wL(v,w))
       =L(v,w)(u,v,w).
\]
In this case the inverse parameter is exactly $[v:w]$.
If $(v,w)=(0,0)$, the point is $[u:0:0]=[1:0:0]$ in the transformed
coordinates, and hence $[x]=[p]$.  Its inverse parameter is defined to
be $[-m:\ell]$.  This pair is nonzero, its $L$-value is zero, and its
$R$-value is $r\ne0$, so it indeed maps to $[1:0:0]$.

It remains to check that the inverse just specified recovers every
parameter.  When $L(s,t)\ne0$, the last two coordinates of $F(s,t)$
recover $[sL:tL]=[s:t]$.  When $L(s,t)=0$, the nonzero vector $(s,t)$
is proportional to $(-m,\ell)$.  For completeness, if $\ell\ne0$,
the equation $\ell s=-mt$ gives
$(s,t)=(t/\ell)(-m,\ell)$ and forces $t\ne0$; if $\ell=0$, then
$m\ne0$, $t=0$, and both projective pairs equal $[1:0]$.
Thus this exceptional case also recovers the original parameter.

Every branch tests an integer equality or inequality.  Matrix products,
polynomial values, signed integer arithmetic, gcd, and exact division
are primitive recursive.  The displayed inverse uses integer adjugate
coordinates followed by normalization, so rational division in the
explanatory proportionality check is not an operation required by the
algorithm.  The formulas define total programs by returning a fixed
default value on inputs failing the decidable hypotheses.  The theorem
asserts their correctness on valid inputs.
\end{proof}

The exceptional value has a geometric interpretation that is already
visible in the arithmetic.  At $(U,V,W)=(1,0,0)$ the gradient of the
transformed conic is $(0,\ell,m)$, so its tangent is
$\ell V+mW=0$.  This is exactly the line represented by $[-m:\ell]$.
Thus the tangent is included in the bijection rather than discarded as
an undefined slope.

\subsection{Quadratic height comparison with computable constants}

For a rational projective point, $\Ht$ denotes the maximum absolute
coordinate of its primitive integer representative.  Define the
following positive integers from the coefficients and the seed:
\begin{equation}\label{eq:conic-height-constants}
\begin{split}
 D_s&=m^2+|bm-c\ell|+|cm|,\\
 D_t&=\ell^2+|a\ell|+|b\ell-am|,\\
 D&=\max(D_s,D_t),\qquad A=3CP,\qquad K=2PD.
\end{split}
\end{equation}
Here $D\geq1$ because $(\ell,m)\ne(0,0)$.

\begin{theorem}[Height of a conic parametrization]
\label{thm:conic-height}
For every $T\in\PP^1(\Q)$,
\begin{equation}\label{eq:conic-height}
 \Ht(T)^2\leq K\Ht(\Phi(T)),
 \qquad
 \Ht(\Phi(T))\leq A\Ht(T)^2.
\end{equation}
In particular, a complete list of conic points of height at most $B$
is obtained by evaluating $\Phi$ at all canonical parameters of height
at most $\lfloor\sqrt{KB}\rfloor$ and retaining exactly those images
of height at most $B$.
\end{theorem}

\begin{proof}
Take a primitive integer pair $(s,t)$ representing $T$, and put
$S=\max(|s|,|t|)\geq1$.
The coefficient formulas give
\[
 |R(s,t)|\leq CS^2,
 \qquad |\ell|+|m|\leq2CP,
 \qquad |L(s,t)|\leq2CPS.
\]
The coordinates of $MF$ are
\[
 -p_0R,\qquad -p_1R+sL,\qquad -p_2R+tL.
\]
Their absolute values are at most $3CPS^2=AS^2$.  Primitive
normalization cannot increase that maximum.  This proves the upper
inequality.

For the lower inequality it is essential to control cancellation of
common factors.  Let
\[
 g=\gcd\bigl(|R(s,t)|,|sL(s,t)|,|tL(s,t)|\bigr).
\]
This integer is positive because $F\ne0$.
Since $\gcd(s,t)=1$, a B\'ezout identity for $s,t$ shows
$\gcd(sL,tL)=|L|$, including the case $L=0$.  Consequently
$g=\gcd(|R|,|L|)$.  The identities
\eqref{eq:conic-identity-s}--\eqref{eq:conic-identity-t} show that
$g$ divides $rs^2$ and $rt^2$.  Since $s^2,t^2$ are coprime, another
B\'ezout identity yields $g\mid r$.  Hence $g\leq|r|$.

Put $T_0=\max(|R|,S|L|)=\max_i|F_i|$.
If $|s|=S$, \eqref{eq:conic-identity-s} gives
\[
 |r|S^2\leq m^2|R|+
 (|bm-c\ell|+|cm|)S|L|\leq D_sT_0\leq DT_0.
\]
If $|t|=S$, the same conclusion follows from
\eqref{eq:conic-identity-t} and $D_t$.
Since $T_0=g\Ht([F])$ and $g\leq|r|$, cancellation of the positive
integer $|r|$ gives
\begin{equation}\label{eq:conic-intermediate-height}
 S^2\leq D\Ht([F]).
\end{equation}
Now let $x$ be the primitive representative of $\Phi(T)=[MF]$.
The integer vector $\operatorname{adj}(M)x$ is nonzero and represents
$[F]$.  Formula \eqref{eq:conic-adjugate} bounds each of its
coordinates by $2P\Ht(\Phi(T))$.  Normalizing it and substituting in
\eqref{eq:conic-intermediate-height} yields
$S^2\leq2PD\Ht(\Phi(T))=K\Ht(\Phi(T))$.

For the last assertion, any image of height at most $B$ has inverse
parameter with squared height at most $KB$.  The parameter box is
therefore complete, and its final height filter removes exactly the
unwanted images.  The integer square root is obtained by a bounded
search through $0,\ldots,KB$.
\end{proof}

The theorem also supplies an inverse height estimate in its strongest
useful elementary form:
\[
 \Ht(\Phi^{-1}(X))\leq\lfloor\sqrt{K\Ht(X)}\rfloor.
\]
The linear bound obtained directly from the adjugate is easier, but
would lose the correct quadratic relation between the parameter
height and the conic height.

\subsection{A finite counting theorem and constructive infinitude}

Let $N_q(B)$ be the number of canonical points of $C_q(\Q)$ with
height at most $B$.  This is a primitive recursive finite count.  In
the next theorem, inequalities with rational-looking constants are
stated after clearing denominators.

\begin{lemma}[Many coprime pairs in a square]
\label{lem:conic-coprime-count}
Let $G(N)$ count pairs $(s,t)$ with $1\leq s,t\leq N$ and
$\gcd(s,t)=1$.  For every $N\geq1$,
\[
 4G(N)\geq N^2.
\]
\end{lemma}

\begin{proof}
For $N=1$ the assertion is immediate.  For $N\geq2$, every pair
with gcd greater than one belongs to at least one of the sets
\[
 E_d=\{(s,t):1\leq s,t\leq N,\ d\mid s,\ d\mid t\},
 \qquad 2\leq d\leq N.
\]
The finite union bound and $\#E_d=\lfloor N/d\rfloor^2$ give
\[
 N^2-G(N)\leq\sum_{d=2}^N\lfloor N/d\rfloor^2
 \leq N^2\sum_{d=2}^N\frac1{d^2}.
\]
There is no infinite series in this argument.  For $d\geq3$,
$1/d^2\leq1/(d(d-1))$, and the finite telescoping identity gives
\[
 \sum_{d=2}^N\frac1{d^2}
 \leq\frac14+\sum_{d=3}^N\left(\frac1{d-1}-\frac1d\right)
 =\frac34-\frac1N.
\]
Thus $G(N)\geq N^2/4+N$, which in particular proves the stated
inequality.  Every displayed rational comparison can equivalently
be multiplied by the positive product of its denominators.
\end{proof}

\begin{theorem}[Explicit linear growth of the point count]
\label{thm:conic-count}
With $A,K$ from \eqref{eq:conic-height-constants},
\[
 N_q(B)\leq4KB\quad(B\geq1),
 \qquad
 B\leq16A N_q(B)\quad(B\geq A).
\]
For each $n\geq1$ one can also produce $n+1$ distinct points on
the conic, each of height at most $An^2$.
\end{theorem}

\begin{proof}
For the upper bound put $S=\lfloor\sqrt{KB}\rfloor\geq1$.
The inverse of every point counted by $N_q(B)$ is a parameter of
height at most $S$.  The number of all nonzero integer pairs in the
box $[-S,S]^2$ is $4S^2+4S$.  Exactly one element in each pair
$\{v,-v\}$ has its first nonzero coordinate positive.  Thus there
are at most $2S^2+2S\leq4S^2$ canonical primitive pairs in this box.
The bijection and $S^2\leq KB$ imply $N_q(B)\leq4KB$.

For the lower bound choose the largest integer $N$ with $AN^2\leq B$.
The condition $B\geq A$ ensures $N\geq1$.
Every positive coprime pair in $[1,N]^2$ gives a different
projective parameter: two primitive positive integer pairs that
are rational multiples of one another are equal.  Their images
are distinct by Theorem~\ref{thm:conic-param} and have height at
most $AN^2\leq B$.  Hence
\[
 4N_q(B)\geq4G(N)\geq N^2.
\]
Maximality of $N$ gives
$B<A(N+1)^2\leq4AN^2\leq16AN_q(B)$, proving the assertion.

For the last statement use the pairwise distinct parameters
$[k:1]$ for $k=0,1,\ldots,n$.  Their heights are at most $n$ when
$n\geq1$, and Theorem~\ref{thm:conic-height} supplies the stated
bound.  This yields an explicit list for each $n$ and therefore
the constructive content of infinitude.
\end{proof}

This is a height-counting result rather than an appeal to an
asymptotic counting theorem.  Both constants and the enumerating
algorithm are given by fixed primitive recursive expressions.
The lower estimate already captures the linear order of growth
of the number of rational points on a pointed conic with respect
to its projective height.

\subsection{Two fully arithmetic examples}

\paragraph{The unit circle.}
For $X^2+Y^2=Z^2$, the preceding construction has the familiar form
\begin{equation}\label{eq:circle-param}
 [s:t]\longmapsto[s^2-t^2:2st:s^2+t^2].
\end{equation}
Its defining identity is
$(s^2-t^2)^2+(2st)^2=(s^2+t^2)^2$.
An inverse is $[Z+X:Y]$ unless $(X:Y:Z)=(-1:0:1)$; that point has
inverse $[0:1]$.  Indeed if $Z+X=0$, the conic equation gives
$Y^2=0$ and hence $Y=0$, so the exceptional point is the only
undefined value of the displayed inverse pair.  Otherwise
substitution into \eqref{eq:circle-param}, followed by use of
$X^2+Y^2=Z^2$, recovers $[X:Y:Z]$.

If $\gcd(s,t)=1$, the common divisor of the three coordinates in
\eqref{eq:circle-param} is exactly
\[
 \epsilon=\begin{cases}2,&s,t\text{ both odd},\\1,&s,t\text{ of opposite parity}.
 \end{cases}
\]
To see this, any common divisor divides $2s^2$ and $2t^2$, and
$\gcd(s^2,t^2)=1$ therefore makes it divide $2$.  If the inputs
are both odd, all three coordinates are even; if their parities
differ, the first and third are odd.  They cannot both be even
because they are coprime.  Since
$|s^2-t^2|\leq s^2+t^2$ and $|2st|\leq s^2+t^2$, we obtain the
exact height formula
\begin{equation}\label{eq:circle-exact-height}
 \Ht(\Phi([s:t]))=\frac{s^2+t^2}{\epsilon}.
\end{equation}
For example $[2:1]$ yields $[3:4:5]$, while $[1:1]$ yields
$[0:1:1]$ after dividing the displayed coordinates by $2$.
The same identity is the elementary connection between rational
points on a conic and primitive Pythagorean triples.

\paragraph{An empty conic detected by a prime-power obstruction.}
The nonsingular conic
\[
 X^2+Y^2=3Z^2
\]
has no rational projective point.  Suppose it had one and choose
a primitive integer representative $(x,y,z)$.  The only square
residues modulo $3$ are $0$ and $1$, so $x^2+y^2\equiv0\pmod3$
forces $3\mid x$ and $3\mid y$.  Write $x=3u,y=3v$ and cancel
$3$ from the equation to obtain $z^2=3(u^2+v^2)$.  The same
square-residue table gives $3\mid z$, contradicting primitivity.
Every step is a finite divisibility calculation.

There \emph{are} projective points of this equation over the
field with three elements, for example $(0:0:1)$.  The obstruction
is the failure of such a residue class to lift to a solution modulo
$9$ with at least one coordinate not divisible by $3$.  Equivalently,
it is the two successive modulo-$3$ deductions just given.  This
distinction matters when formulating a correct finite local
certificate.

\subsection{Why these arguments fit the weak constructive theory}

No existence claim above asks for an unbounded search.  A seed point
is part of the hypotheses of the parametrization theorem; the
theorem does not search indefinitely for one.  A conic with no
supplied seed still has the bounded point enumerator, and a
specific local obstruction may certify emptiness as in the second
example.  A general local-to-global theorem for conics is a further
theorem, not an unstated step of this chapter.

The algebraic parts use fixed polynomial identities, gcd and
B\'ezout algorithms, and decidable branches.  The counting parts
only use bounded ranges.  Their induction invariants are equalities
and inequalities of primitive recursive counting functions; a
bounded finite-union assertion is represented by its primitive
recursive characteristic function.  Finite telescoping is proved
by quantifier-free induction, after clearing the positive
denominators or maintaining a rational pair code.  Distinctness of
the produced list is also a bounded predicate.

Thus the internal assertion of constructive infinitude is the
explicit arithmetic statement that for every input $n$, a specified
primitive recursive list has length $n+1$, distinct valid entries,
and the displayed height bound.  It does not quantify over an
infinite set, invoke a choice principle, or require induction on
an arbitrary formula.  The bijection and the counting estimates
are uniform first-order arithmetic theorems about the supplied
coefficient and point codes.

\section{Finite descent and a complete elliptic-curve example}
\label{sec:descent}

Height-bounded enumeration becomes much stronger when a geometric transformation
strictly decreases an integer measure.  We now give an example in which descent
settles the entire rational-point problem, with no height cutoff left over.  The
main conclusions are Fermat's square-area theorem and the complete list of
rational points on
\[
  E_1:\qquad Y^2Z=X^3-XZ^2.
\]
The argument is elementary, but its logical implementation requires care:
``choose a least counterexample'' is replaced by a specified number of
primitive-recursive steps.  Fermat's result and its connection with the
congruent-number curves are classical; see \cite{ConradCongruent}.  An existing
Coq treatment of Fermat's theorem is useful background for formalization
\cite{DelahayeMayero}.  Neither reference is used as an unproved arithmetic
premise below.

\subsection{A descent principle with a numerical stopping bound}

\begin{lemma}[Finite descent in \iPRA]\label{lem:finite-descent}
Fix primitive-recursive functions $D$ and $\rho$ on codes of states, and a
primitive-recursive Boolean predicate $G$.  Suppose that \iPRA\ proves
\[
 G(s)\ \Longrightarrow\
  \bigl(0<\rho(s)\ \wedge\ G(D(s))\ \wedge\ \rho(D(s))<\rho(s)\bigr).
 \tag{D}
\]
Then \iPRA\ proves $\neg G(s)$ for every state $s$.
\end{lemma}

\begin{proof}
Define $s_0=s$ and $s_{j+1}=D(s_j)$ by primitive recursion.  Under the
assumption $G(s)$, quantifier-free induction on $j$ gives
\[
 G(s_j)\quad\hbox{and}\quad \rho(s_j)+j\leq\rho(s).
 \tag{1}
\]
Indeed, (D) gives $G(s_{j+1})$ and
$\rho(s_{j+1})+1\leq\rho(s_j)$, which proves the induction step by addition.
At $j=\rho(s)$, (1) implies $\rho(s_j)=0$, while (D) and $G(s_j)$ imply
$0<\rho(s_j)$.  This is a contradiction obtained from the assumption $G(s)$,
and therefore proves its negation intuitionistically.  There is no use of
double-negation elimination.  The displayed induction predicate is
quantifier-free after replacing the Boolean tests and inequalities by their
primitive-recursive characteristic functions.
\end{proof}

This is a schema for fixed primitive-recursive functions with their defining
equations.  It does not posit an internal universal evaluator for arbitrary
primitive-recursive programs.  In the application, the number of descent steps
is the original hypotenuse.  An abstract assertion that there is some smaller
counterexample would not by itself supply the required function $D$; the
formulas below do.

\subsection{Arithmetic tools without prime factorization}

All integers in the next three lemmas are nonnegative unless a sign is shown.
Division is used only after divisibility has been established.  Each exact
quotient can equivalently be implemented by a bounded search.

\begin{lemma}[The gcd tools]\label{lem:descent-gcd}
The following operations and implications have primitive-recursive witnesses
and quantifier-free-induction proofs.
\begin{enumerate}
\item If $a,b>0$, the greatest common divisor $g$ is a positive common divisor,
and there are signed integers $u,v$ with $ua+vb=g$.
\item If $\gcd(a,b)=1$ and $a\mid bc$, then $a\mid c$.
\item If $d^2\mid c^2$ and $d>0$, then $d\mid c$.
\item If $a,b,c>0$, $\gcd(a,b)=1$, and $ab=c^2$, then one can compute
positive integers $r,s$ such that $a=r^2$, $b=s^2$, and $c=rs$.
\end{enumerate}
\end{lemma}

\begin{proof}
For the first assertion, order the initial pair as
$(\max(a,b),\min(a,b))$ and perform Euclidean division.  Replace a pair
$(u,v)$ with $v>0$ by $(v,u\bmod v)$; after reaching a zero second entry,
leave the pair unchanged.  Before stopping, the sum of the entries strictly
decreases.  Thus $a+b+1$ iterations suffice.  Track the two entries as signed
linear combinations of the original $a,b$: subtraction of the quotient times
the second combination from the first gives the next combination.  The common
divisors are unchanged at each step, and the last nonzero entry divides both
preceding entries, hence both original entries.  The tracked last combination
is the required B\'ezout identity.  These statements are invariants of a fixed
finite computation and use quantifier-free induction.  If an explicit bound
for the coefficient storage is desired, each quotient is at most
$M=\max(a,b)$ and each coefficient is bounded in absolute value by
$(M+1)^{a+b+2}$.

For the second assertion, multiply a B\'ezout identity $ua+vb=1$ by $c$.
Both terms on the resulting right side are divisible by $a$ whenever
$a\mid bc$.  The same identity also shows that coprimality is preserved when
one of the entries is squared: squaring $ua+vb=1$ gives a B\'ezout identity
between $a$ and $b^2$.

For the third assertion, $c=0$ gives $d\mid c$ immediately. Otherwise
put $h=\gcd(c,d)$, $c=hr$, and $d=hs$.
Then $\gcd(r,s)=1$, and cancellation in $d^2\mid c^2$ gives $s^2\mid r^2$.
In particular $s\mid r^2$.  Applying the preceding coprimality observation
and the second assertion yields $s\mid1$, so $s=1$.  Consequently $d=h$ divides
$c$.

For the fourth assertion, put $h=\gcd(a,c)$, $a=hu$, and $c=hv$.
Here $\gcd(u,v)=1$, and the equation becomes
\[
 ub=hv^2.
\]
Since $u$ is coprime to $v^2$, it divides $h$.  Write $h=ue$.  Cancellation
now gives $a=e u^2$ and $b=e v^2$.  Thus $e$ divides both $a$ and $b$;
coprimality forces $e=1$.  Taking $r=u$ and $s=v$ proves the claim, including
$c=rs$.  This proof extracts the square roots by gcd and exact division.
It does not need a prime factorization of the inputs.
\end{proof}

Repeatedly applying the last assertion splits a square product of any fixed
number of pairwise coprime positive factors into squares.  To justify this
iteration, if $a$ is coprime to each of $b$ and $c$, multiply their B\'ezout
identities to obtain a B\'ezout identity for $a$ and $bc$.  Thus one factor is
coprime to the product of all the remaining factors.

\begin{lemma}[Normalization of square-area triangles]\label{lem:triangle-normalize}
From positive integers $a,b,c,s$ satisfying
\[
 a^2+b^2=c^2,\qquad ab=2s^2,
 \tag{2}
\]
one can compute positive integers $A,B,C,S$ satisfying the same equations
and $\gcd(A,B)=1$, with $C\leq c$.
\end{lemma}

\begin{proof}
Let $g=\gcd(a,b)$ and put $A=a/g$, $B=b/g$.  Since
$c^2=g^2(A^2+B^2)$, Lemma~\ref{lem:descent-gcd}(3) gives $g\mid c$;
put $C=c/g$.  The resulting Pythagorean equation is primitive in the two
legs.  Its legs cannot both be even, by coprimality, and cannot both be odd,
because their squared sum would be $2$ modulo $4$.  Thus exactly one leg is
even, and $K=AB/2$ is an integer.  Equation (2) now says $s^2=g^2K$.
Applying the same square-divisibility lemma gives $g\mid s$.  Put $S=s/g$.
Cancelling $g^2$ in the two equations proves the result.  All four outputs
are positive, and $C=c/g\leq c$.
\end{proof}

\begin{lemma}[Primitive Pythagorean parametrization]\label{lem:pythagorean-descent}
Let $A,B,C>0$, $A^2+B^2=C^2$, and $\gcd(A,B)=1$.  After exchanging the legs
so that $A$ is even, there are unique positive integers $m>n$ such that
\[
 A=2mn,\qquad B=m^2-n^2,\qquad C=m^2+n^2,
 \tag{3}
\]
where $\gcd(m,n)=1$ and $m,n$ have opposite parity.  They are obtained by
bounded arithmetic from $A,B,C$.
\end{lemma}

\begin{proof}
The parity argument in the preceding lemma shows that $B$ and $C$ are odd.
Also $\gcd(B,C)=1$: any common divisor $d$ satisfies $d^2\mid A^2$, hence
$d\mid A$ by Lemma~\ref{lem:descent-gcd}(3), and then $d=1$.
As $C>B>0$, the integers
\[
 U=\frac{C+B}{2},\qquad V=\frac{C-B}{2}
\]
are positive, with $U>V$.  A common divisor divides their sum $C$ and their
difference $B$, so $\gcd(U,V)=1$.  Their product is
\[
 UV=\frac{C^2-B^2}{4}=\left(\frac A2\right)^2.
\]
The square-splitting lemma gives $U=m^2$, $V=n^2$, and $A/2=mn$.
Addition and subtraction give (3), while $U>V$ gives $m>n$.
A common divisor of $m,n$ would divide $U,V$, so they are coprime.
Their squares sum to the odd number $C$, so they have opposite parity.
Finally, $m^2=(C+B)/2$ and $n^2=(C-B)/2$ determine the positive roots
uniquely: $x^2=y^2$ for nonnegative integers implies $x=y$ by
$(x-y)(x+y)=0$, or directly by strict monotonicity of squaring.
Each root is at most $C$ and can be found by bounded search.
\end{proof}

\subsection{The explicit descent}

\begin{theorem}[Fermat's square-area theorem in \iPRA]\label{thm:fermat-area}
There are no positive integers $a,b,c,s$ satisfying (2).  Consequently there
is no rational right triangle of positive side lengths whose area is a
nonzero rational square.  In particular, a rational right triangle cannot
have area $1$.
\end{theorem}

\begin{proof}
Normalize a supposed solution by Lemma~\ref{lem:triangle-normalize}.  We
will transform every primitive solution into another primitive solution
with a smaller hypotenuse.

Apply Lemma~\ref{lem:pythagorean-descent} and substitute (3) into the
square-area equation.  Cancelling $2$ gives
\[
 m n(m+n)(m-n)=S^2.                         \tag{4}
\]
These four factors are positive and pairwise coprime.  For pairs involving
$m$ or $n$, a common divisor would also divide both $m$ and $n$.
For the remaining pair, a common divisor of $m+n$ and $m-n$ divides
$2m$ and $2n$, and hence divides $2$ by a B\'ezout identity for $m,n$.
Both $m+n$ and $m-n$ are odd, so that common divisor is $1$.
The square-splitting lemma applied three times to (4) supplies positive
integers $r,s,t,u$ with
\[
 m=r^2,\qquad n=s^2,\qquad m+n=t^2,\qquad m-n=u^2.
 \tag{5}
\]
In particular $t,u$ are odd, $t>u>0$, and $\gcd(t,u)=1$.
Moreover $s$ is even.  Indeed $m,n$ have opposite parity; if $m$ were even
and $n$ odd, then the first two equations of (5) would give
$m-n\equiv-1\equiv3\pmod4$, contradicting $m-n=u^2$.
Thus $n$ is even, and $n=s^2$ forces $s$ even.

Define
\[
 A'=\frac{t+u}{2},\qquad B'=\frac{t-u}{2},\qquad
 C'=r,\qquad S'=\frac s2.                  \tag{6}
\]
All four are positive integers.  Direct calculation using (5) gives
\begin{align*}
 (A')^2+(B')^2
   &=\frac{t^2+u^2}{2}=m=r^2=(C')^2,\\
 A'B'&=\frac{t^2-u^2}{4}=\frac n2
       =\frac{s^2}{2}=2(S')^2.
\end{align*}
A common divisor of $A',B'$ divides their sum $t$ and difference $u$,
so the new triangle is primitive.  Finally,
\[
 0<C'=r\leq r^2=m<m^2+n^2=C.
 \tag{7}
\]
This completes the explicit descent step.

For its \iPRA\ implementation, let $G$ be the Boolean predicate that a
quadruple is positive, satisfies (2), and has coprime legs.  On a valid input,
let $D$ carry out the leg exchange, the bounded square-root extractions in
(3) and (5), and formula (6).  On an invalid input let $D$ return the zero
quadruple.  Bounded integer square root, gcd, exact quotient, comparison,
and tuple coding are primitive recursive, so this defines a total
primitive-recursive function.  The preceding paragraphs prove (D) with
$\rho(a,b,c,s)=c$.  Lemma~\ref{lem:finite-descent} therefore proves that
$G$ is always false.  Normalization then rules out (2) without the
coprimality condition.

For the rational statement, suppose positive rationals $a,b,c$ satisfy
$a^2+b^2=c^2$ and $ab/2=q^2$ with $q\ne0$.  Take the product $L>0$ of the
positive denominators of $a,b,c,q$.  The numbers
$La,Lb,Lc,L|q|$ are positive integers satisfying (2), a contradiction.
The product of denominators is a specified primitive-recursive operation;
no least common denominator is needed.
\end{proof}

The proof has a useful complexity audit.  On a valid state of hypotenuse
$C$, the parameter searches need inspect no integer exceeding $C$.
The next state's coordinates are all at most $C$: the two legs are smaller
than the new hypotenuse, and its square-area root is smaller than that
hypotenuse as well.  Iteration for exactly the original $C$ steps is
therefore more than sufficient.  The contradiction concerns the purported
input itself, so it does not require an a priori bound on all possible
counterexamples.

\begin{corollary}[A quartic obstruction]\label{cor:quartic-obstruction}
There are no positive integers $x>y$ and $z$ with $x^4-y^4=z^2$.
In particular $a^4+b^4=c^4$ has no solution in positive integers.
\end{corollary}

\begin{proof}
A solution of the first equation would give the integer right triangle
\[
 \bigl(2x^2y^2,\ x^4-y^4,\ x^4+y^4\bigr).
\]
The identity $(2uv)^2+(u^2-v^2)^2=(u^2+v^2)^2$, with $u=x^2$ and
$v=y^2$, verifies that it is right-angled.  Its area is
\[
 x^2y^2(x^4-y^4)=(xyz)^2>0,
\]
contrary to Theorem~\ref{thm:fermat-area}.  For the second assertion,
$c^4-a^4=(b^2)^2$ would contradict the first, since $c>a>0$.
\end{proof}

\begin{corollary}[Three rational squares]\label{cor:square-progression}
There are no rational numbers $r,s,t,d$, with $d\ne0$, for which
$s^2-r^2=t^2-s^2=d^2$.
\end{corollary}

\begin{proof}
Put $R=|r|$, $S=|s|$, $T=|t|$.  Then $T>R\geq0$ and $S>0$.
The positive rational numbers
\[
  T-R,\quad T+R,\quad 2S
\]
form a right triangle because
$(T-R)^2+(T+R)^2=2(T^2+R^2)=4S^2$.
Its area is $(T^2-R^2)/2=d^2>0$, again a contradiction.
\end{proof}

\subsection{A rational correspondence with cubic curves}

For a positive rational $n$, write
\[
 E_n:\quad y^2=x^3-n^2x.
\]
Its projective equation is $Y^2Z=X^3-n^2XZ^2$; after clearing the denominator
of $n$, this is an integral homogeneous equation in the format used earlier.
We first keep signs in the triangle coordinates.  This makes the formulas
mutually inverse without dividing into geometric sign cases.

\begin{theorem}[Congruent-number correspondence]\label{thm:En-correspondence}
For each positive rational $n$, there are mutually inverse rational maps
between the following two collections of rational tuples:
\[
 \{(a,b,c):a^2+b^2=c^2,\ ab=2n\}
 \quad\longleftrightarrow\quad
 \{(x,y):y^2=x^3-n^2x,\ y\ne0\}.
\]
They are
\begin{align}
 (a,b,c)&\longmapsto
   \left(\frac{nb}{c-a},\frac{2n^2}{c-a}\right),
   \label{eq:triangle-to-curve}\\
 (x,y)&\longmapsto
   \left(\frac{x^2-n^2}{y},\frac{2nx}{y},
                     \frac{x^2+n^2}{y}\right).
   \label{eq:curve-to-triangle}
\end{align}
Every indicated denominator is nonzero.  These are primitive-recursive maps
on reduced rational codes, with any fixed default output outside their stated
domains.  Thus $E_n$ has a rational point with $y\ne0$ if and only if there
is a positive rational right triangle of area $n$.
\end{theorem}

\begin{proof}
If $(a,b,c)$ is in the first collection, then $ab=2n\ne0$, so $a,b$ are
nonzero.  Also $c-a\ne0$: otherwise $b^2=c^2-a^2=0$ and hence $b=0$.
Set $\delta=c-a$, and define $x,y$ by
\eqref{eq:triangle-to-curve}.  Then $y\ne0$.  The two identities
\[
 b^2-\delta^2=2a\delta,\qquad
 b^2+\delta^2=2c\delta
\]
follow by expanding and using $a^2+b^2=c^2$.  They imply
\[
 x^2-n^2=ay,\qquad x^2+n^2=cy,
 \qquad 2nx=by.
 \tag{8}
\]
In addition $ax=anb/\delta=2n^2/\delta=y$.  Multiplying the first
identity in (8) by $x$ gives
$x(x^2-n^2)=axy=y^2$, so the image lies on $E_n$.
The three identities in (8) show that applying
\eqref{eq:curve-to-triangle} recovers exactly $a,b,c$.

Conversely, let $(x,y)$ lie on $E_n$ with $y\ne0$.  Since
$y^2=x(x^2-n^2)$, both $x$ and $x^2-n^2$ are nonzero.  Define $a,b,c$
by \eqref{eq:curve-to-triangle}.  The identity
\[
 (x^2-n^2)^2+(2nx)^2=(x^2+n^2)^2
\]
gives $a^2+b^2=c^2$, and
\[
 ab=\frac{2nx(x^2-n^2)}{y^2}=2n.
\]
Furthermore $c-a=2n^2/y\ne0$, so substitution into
\eqref{eq:triangle-to-curve} returns $x,y$.  This proves both the domain
checks and the inverse assertions.

For positivity, $ab=2n>0$ makes $a,b$ nonzero with the same sign; also
$c^2=a^2+b^2>0$.  Replacing $a,b,c$ by their absolute values therefore
produces positive lengths and preserves both equations.  Conversely a
positive triangle is already a tuple in the first collection.  All sign
tests and rational operations are decidable and primitive recursive on
integer numerator and positive denominator codes.
\end{proof}

As an exact numerical check, the triangle $(3,4,5)$ has area $6$ and gives
$(x,y)=(12,36)$ on $E_6$, since $36^2=12^3-36\cdot12$.  Formula
\eqref{eq:curve-to-triangle} sends it back to $(3,4,5)$.

\begin{theorem}[Complete rational points of $E_1$]\label{thm:E1}
In the coding of rational projective points, \iPRA\ proves that
\[
 E_1(\Q)=\{[0:1:0],\ [0:0:1],\ [1:0:1],\ [-1:0:1]\}.
 \tag{9}
\]
The braces in (9) abbreviate the finite disjunction that every rational
projective point satisfying the cubic equation is equal to one of the four
listed points, together with the direct verification that all four are on
the cubic.
Projective notation permits either sign; the last displayed point has
canonical integer code $(1,0,-1)$ under our first-nonzero-positive convention.
\end{theorem}

\begin{proof}
Let $[X:Y:Z]$ satisfy $Y^2Z=X^3-XZ^2$.
The equality $Z=0$ is decidable.  In that case $X^3=0$, so $X=0$, and
nonzeroness of the coordinate tuple gives $Y\ne0$.  The point is therefore
$[0:1:0]$.

If $Z\ne0$, put $x=X/Z$ and $y=Y/Z$.  They satisfy
$y^2=x^3-x$.  The equality $y=0$ is again decidable.  If $y\ne0$,
Theorem~\ref{thm:En-correspondence} produces a positive rational right
triangle of area $1$, contradicting Theorem~\ref{thm:fermat-area}.
Thus $y=0$, and $x(x-1)(x+1)=0$.  Rational equality is decidable and
$\Q$ has no zero divisors: test $x=0$, and if it is false cancel $x$;
then test $x-1=0$, and if it is false cancel $x-1$.  This yields
$x=0$, $x=1$, or $x=-1$.  They give exactly the three affine points in
(9).  Substitution verifies each of the four displayed projective points.
\end{proof}

For the usual geometric interpretation, the cubic is nonsingular.  With
$F=Y^2Z-X^3+n^2XZ^2$, its partial derivatives are
\[
 F_X=-3X^2+n^2Z^2,\qquad F_Y=2YZ,\qquad
 F_Z=Y^2+2n^2XZ.
\]
At the unique point at infinity, $F_Z\ne0$.  In the affine chart, simultaneous
vanishing of $F$ and $F_Y$ gives $y=0$ and $x\in\{0,n,-n\}$.
At these points $F_X$ is respectively $n^2,-2n^2,-2n^2$, all nonzero.
The calculation is valid over every characteristic-zero field with $n\ne0$.
It identifies $E_n$ as a nonsingular plane cubic with a rational distinguished
point, the customary elliptic-curve situation.  The point classification
itself is entirely the integer and rational calculation above; it requires
no group-law or general finite-generation theorem.

\subsection{What the formalization must establish}

The proof naturally separates into five reusable Lean components:
\begin{enumerate}
\item The signed-integer and rational-code algebra, including bounded gcd,
exact division, and the four implications in
Lemma~\ref{lem:descent-gcd}.
\item Primitive Pythagorean normalization and the parameter extraction
of Lemmas~\ref{lem:triangle-normalize} and
\ref{lem:pythagorean-descent}.
\item A total function \texttt{descendSquareArea} on quadruple codes, with
its preservation and strict-decrease theorems proved under the Boolean
validity hypothesis.
\item A general finite-iteration lemma using only the displayed
quantifier-free invariant, followed by
\texttt{noSquareAreaTriangle}.
\item The two rational transformations, denominator proofs, inverse
identities, and the finite disjunction \texttt{rationalPointsE1}.
\end{enumerate}

The included executable reference file \texttt{code/descent.py} implements
the arithmetic transformations and checks their identities on exact rational
examples.  It also implements the total descent map.  A valid input for its
main descent branch cannot exist, by the theorem; therefore finite testing
cannot exercise that branch on an actual valid counterexample or replace
its mathematical proof.  The branch's preservation proof is the symbolic
calculation (4)--(7).  The remaining formalization task is to check those
symbolic calculations and their arithmetic preliminaries in Lean, and then
to connect the fixed algorithms and the indicated quantifier-free
inductions to the chosen object-theory representation of \iPRA.

\section{Local obstructions and denominators on cubic curves}
\label{sec:local}

The previous sections give both positive descriptions and a global negative
result. Two further tools make the interface between geometry and finite
arithmetic explicit: reduction modulo a supplied integer, and the special
denominator structure of points on an integral Weierstrass equation.

\subsection{Finite congruence certificates}

For $m\ge2$, call a residue vector
$\bar x=(\bar x_0,\ldots,\bar x_n)\in\{0,\ldots,m-1\}^{n+1}$
\emph{unimodular modulo $m$} when
$\gcd(m,\bar x_0,\ldots,\bar x_n)=1$. For prime $m$, this means exactly
that not all coordinates are zero. For composite $m$, nonzeroness alone
is insufficient.

\begin{theorem}[A checkable certificate of no rational point]
\label{thm:modular}
Let $F_1,\ldots,F_s$ be homogeneous integer polynomials. If none of the
unimodular residue vectors modulo $m$ satisfies all equations
$F_i(\bar x)\equiv0\pmod m$, their projective zero locus has no rational
point. The premise is a primitive recursive test of at most $m^{n+1}$
vectors, and $\iPRA$ proves the implication.
\end{theorem}
\begin{proof}
A rational point has a primitive integral representative $x$. Reduce every
coordinate to $\bar x_i$ in $[0,m-1]$. Any common divisor of $m$ and all
$\bar x_i$ divides every $x_i$, so it is one; hence the residue vector is
unimodular. Polynomial evaluation respects congruence: each product does,
by induction on its number of factors, and a finite sum does as well.
Thus $F_i(x)=0$ implies $F_i(\bar x)\equiv0\pmod m$, contradicting the
tested premise. This is an intuitionistic proof of a negation from an
assumed point, not an appeal to a converse local--global principle.
\end{proof}

The empty conic $X^2+Y^2=3Z^2$ of Section~\ref{sec:conics} illustrates
why the modulus matters. Modulo three it has the nonzero vectors
$(0,0,1)$ and $(0,0,2)$. Modulo nine it has no unimodular solution: reduction
modulo three forces $3\mid X,Y$, and the original congruence modulo nine
then forces $3\mid Z$. The exact companion reproduces both residue counts.
Testing just the affine chart $Z=1$ would not be a valid general method
for excluding projective points, since a rational point may reduce to a
different chart.

The theorem is a one-way certificate. Failure to find an obstruction among
the tested moduli supplies no global existence conclusion. There is no
unbounded search for a successful modulus in any result here.

\subsection{Denominators on integral Weierstrass equations}

Consider
\[
      E_{A,B}:\quad y^2=x^3+Ax+B,\qquad A,B\in\Z.
\]
If $4A^3+27B^2\ne0$, its projective closure
$Y^2Z=X^3+AXZ^2+BZ^3$ is nonsingular. This can be checked directly.
At infinity the equation forces $X=0$, and the derivative with respect to
$Z$ at $[0:1:0]$ is $1$. At an affine singular point, the derivatives give
$y=0$ and $3x^2+A=0$, while the equation gives $x^3+Ax+B=0$.
Elimination gives $A=-3x^2$, $B=2x^3$ and therefore
$4A^3+27B^2=0$, contrary to the hypothesis. These arguments are polynomial
identities over characteristic zero. We need no group law for the results
of this section. Indeed their denominator assertions remain valid even
when the cubic is singular.

\begin{theorem}[Square and cube denominators]\label{thm:denominators}
Every rational affine point on $E_{A,B}$ has a unique expression
\[
        x=\frac a{d^2},\qquad y=\frac b{d^3},\qquad
        d>0,\quad\gcd(a,d)=\gcd(b,d)=1,
\]
with $a,b\in\Z$. These integers satisfy
\begin{equation}\label{eq:elliptic-integral}
             b^2=a^3+Aad^4+Bd^6.
\end{equation}
Conversely, integers with these conditions give a rational point. The
conversion in both directions is primitive recursive.
\end{theorem}
\begin{proof}
Write $x=a/q$ and $y=c/e$ in reduced form, with $q,e>0$. Let
\[
             T=a^3+Aaq^2+Bq^3.
\]
Then $y^2=T/q^3$. Since $T\equiv a^3\pmod q$ and $\gcd(a,q)=1$,
we have $\gcd(T,q)=1$, hence $\gcd(T,q^3)=1$. Also
$c^2/e^2$ is a reduced fraction. Uniqueness of reduced fractions therefore
gives
\[
                         e^2=q^3,\qquad c^2=T.
\]
For clarity, the uniqueness assertion follows from cross multiplication
and Bezout: $u/v=u'/v'$ with both denominators positive and both fractions
coprime implies $v\mid v'$ and $v'\mid v$, then $v=v'$ and $u=u'$.
It remains valid when a numerator is zero, since its reduced denominator
must then be one.

In the factorization of $q,e$, every prime exponent satisfies $2v=3u$.
Thus $u$ is even and $v$ is three times its half. The finite factorization
of Lemma~\ref{lem:arithmetic} constructs $d>0$ with $q=d^2$, $e=d^3$.
Equivalently the algorithm may take $d=\lfloor\sqrt q\rfloor$, since the
proof has established that $q$ is a square. Put $b=c$. Coprimality and
\eqref{eq:elliptic-integral} now follow from the reduced fractions.
The converse is obtained by dividing the integer equation by $d^6$.
For uniqueness, the reduced denominator of $x$ is $d^2$, so $d$ is its
unique positive square root and both numerators are determined.
\end{proof}

\begin{theorem}[Complete elliptic search at a given $x$-height]
\label{thm:elliptic-search}
Let $T\ge1$. All rational affine points on $E_{A,B}$ with
$\Ht(x)\le T$ are obtained by the following finite procedure:
\begin{enumerate}
\item Visit $1\le d\le\lfloor\sqrt T\rfloor$ and $-T\le a\le T$
with $\gcd(a,d)=1$.
\item Compute $W=a^3+Aad^4+Bd^6$. If $W\ge0$, compute its integer
square root $b_0$. Retain $(a/d^2,\pm b_0/d^3)$ exactly when $b_0^2=W$,
counting the two signs only once if $b_0=0$.
\end{enumerate}
The list is duplicate-free and has at most
$2(2T+1)\lfloor\sqrt T\rfloor$ points. Adjoining $[0:1:0]$ accounts
for the only point at infinity. This construction and its correctness are
provable in $\iPRA$.
\end{theorem}
\begin{proof}
Every retained pair satisfies the equation by direct substitution.
Coprimality of $a,d$ implies that $a/d^2$ is reduced. It also forces
$\gcd(b_0,d)=1$: a common prime would divide $a^3$ after reducing
\eqref{eq:elliptic-integral} modulo that prime, and hence would divide
$a$, a contradiction. Thus the displayed pairs are canonical, and the
uniqueness in the denominator theorem excludes duplicates beyond the
explicit zero-sign case.

Conversely the denominator theorem puts every point of the required height
in the form $a/d^2,b/d^3$ with $|a|\le T$ and $d^2\le T$. That pair
$(a,d)$ is visited and its value $W=b^2$ passes the square test. The sign
of $b$ is included. This proves completeness and the count bound.
For an explicit bound on the square-root search,
\[
      0\le W\le(1+|A|+|B|)T^3
\]
whenever $W\ge0$. If $C=1+|A|+|B|$, its square root is at most $CT^2$,
since $CT^3\le C^2T^4$. All loops and stored integers are therefore
explicitly bounded. At infinity, $Z=0$ forces $X=0$ and leaves the single
projective point $[0:1:0]$.
\end{proof}

This makes the familiar denominator geometry an actual bounded algorithm.
It does not provide a global height bound for all points on an arbitrary
cubic. For $E_1$, that global problem was already settled by the descent
proof: the affine list is exactly $(-1,0),(0,0),(1,0)$. For
$y^2=x^3-36x$, the point $(12,36)$ gives the rational right triangle
$(3,4,5)$ through the formulas in Section~\ref{sec:descent}. Thus the
same equations support both obstructed and solvable examples.

\section{A constructive interface for height descent and finite generation}
\label{sec:generation}

An important role of height in Diophantine geometry is to reduce all points
to finitely many small ones. The following theorem isolates a version of
that argument in which the required witnesses are supplied by actual
primitive recursive functions. It is a theorem about explicitly encoded
groups, not an assertion that every elliptic curve already comes equipped
with these data.

\begin{theorem}[Finite generation from a computable reducing step]
\label{thm:generation}
Fix an abelian group represented by natural-number codes, with primitive
recursive operations, equality and a validity predicate. Suppose $\iPRA$
proves its group identities on valid codes. Let $h$ be a primitive recursive
natural-valued height. Fix $m\ge2$, a threshold $T$, and finite lists of valid group elements
\[
      R=(r_0,\ldots,r_{s-1}),\qquad S=(s_0,\ldots,s_{t-1}).
\]
Suppose the following data and their correctness proofs are given:
\begin{enumerate}
\item $S$ lists all valid group elements of height at most $T$;
\item primitive recursive functions $D(P)$ and $i(P)$ such that, for every
valid $P$ with $h(P)>T$, the point $D(P)$ is valid, $i(P)<s$, and
\[
               P=mD(P)+r_{i(P)},\qquad h(D(P))<h(P).
\]
\end{enumerate}
Then $R\cup S$ generates the group. More precisely, a primitive recursive
algorithm produces, for each valid $P$, explicit integer coefficients
expressing $P$ in these generators. Its correctness is provable in $\iPRA$.
\end{theorem}
\begin{proof}
Keep a current point, an active-step count $k$, and a coefficient vector
$(a_0,\ldots,a_{s-1})$. Start at $P$, $k=0$, and all coefficients zero.
When the current point $Q$ has height at most $T$, keep the whole state
unchanged. Otherwise first save $i=i(Q)$ and $D(Q)$ from the old point,
then replace $Q$ by $D(Q)$, add $m^k$ to $a_i$, and increase $k$ by one.
Run this stopped iteration for exactly
$h(P)\mathbin{\dotminus}T$ stages, where $\dotminus$ is truncated
subtraction.

The rank $h(Q)\mathbin{\dotminus}T$ is zero precisely in the stopped
case. In an active step it decreases by at least one; in a stopped step
it remains zero. Quantifier-free induction therefore proves that after
$j$ stages it is at most
$(h(P)\mathbin{\dotminus}T)\mathbin{\dotminus}j$. The final point has
height at most $T$ and hence is some $s_\ell$ in the supplied list. Its
index is found by bounded equality search.

Throughout the iteration the following group identity is maintained:
\begin{equation}\label{eq:descent-expansion}
                   P=m^k Q+\sum_{i<s}a_i r_i.
\end{equation}
It holds initially. An active step substitutes
$Q=mD(Q)+r_{i(Q)}$ into the right side and uses distributivity of
integer multiplication in an abelian group. The result is exactly the
coefficient update. A stopped step changes nothing. At the final state,
replace $Q$ by $s_\ell$ in \eqref{eq:descent-expansion}; this is the
promised expression. All coefficients can even be chosen nonnegative
with this particular supplied reduction.

The loops have the given numerical fuel. Scalar multiplication is repeated
group addition, so its definition is primitive recursive. Every exponent
is at most $h(P)\dotminus T$, and every coefficient is bounded by
$\sum_{j\le h(P)}m^j$. Equation \eqref{eq:descent-expansion} and the rank
inequality are quantifier-free statements about the coded operations, so
their inductions are instances of the permitted rule. No choices from an
unspecified finite quotient are made.
\end{proof}

A commonly useful sufficient numerical condition for strict decrease is
\[
              2h(D(P))\le h(P)+T\qquad(h(P)>T).
\]
Indeed $h(P)+T<2h(P)$ implies $h(D(P))<h(P)$ by integer arithmetic.
This avoids real logarithms and limit-defined canonical heights in the
interface. Constructing such a height and a reducing selector for a
particular group is an additional mathematical obligation.

\subsection{What this does and does not establish about larger theorems}

Theorem~\ref{thm:generation} captures the final reduction step of a
Mordell--Weil-style proof. A statement that a quotient group is finite is
not, on its own, the supplied primitive recursive selector required here.
One must provide representatives, produce the corresponding preimages,
and prove the height decrease in the chosen arithmetic. Likewise a proof
of finiteness of rational points is not automatically an explicit global
height bound to which Corollary~\ref{cor:heightbound} applies.

The results actually completed in this paper need none of those additional
inputs: pointed conics come with explicit formulas, rational preperiodic
points have the resultant-derived height bound, and the curve $E_1$ is
settled by its concrete decreasing transformation. The broad classical
theorems are valuable targets for further work, but they are not hidden
premises in these arguments. No inference here rests on an assertion that
an unexamined classical proof must formalize in weak arithmetic.

\subsection{Logical inventory of the completed proofs}

The paper uses first-order numerical quantification to state its theorems,
including unbounded universal claims about all rational points. Their
proofs use the following finite mechanisms:
\begin{center}
\begin{tabular}{p{0.31\textwidth}p{0.60\textwidth}}
\toprule
Mathematical operation & Arithmetic implementation \\
\midrule
Projective equivalence & Gcd normalization and Bezout uniqueness \\
Finite point loci & Explicit coordinate boxes and polynomial tests \\
Coordinate morphisms & Integral homogeneous formulas and nonvanishing guards \\
Height lower bounds & Polynomial certificates and bounds on their common divisors \\
Infinite families & An injective formula evaluated on a supplied finite initial segment \\
Preperiodicity & A proved height cutoff and a finite repeated-state search \\
Nonexistence by descent & Iteration for the initial integer rank \\
Local impossibility & Exhaustive residue tests for a supplied modulus \\
Finite generation & A supplied reducing selector and a telescoping group identity \\
\bottomrule
\end{tabular}
\end{center}
Each induction has an equation, a Boolean invariant, or a bounded
arithmetical invariant. There is no quantification over arbitrary sets of
natural numbers, no completeness principle for $\R$, and no appeal to an
unspecified unbounded search. This is the substantive reason the development
fits the chosen constructive arithmetic.

\section{Lean verification and a path to complete internalization}
\label{sec:lean}

\subsection{What has actually been checked}

The archive contains an arithmetic core whose source was successfully
compiled. The file \texttt{lean/DiophantineCore.lean} was checked
with Lean 4.32.0, commit
\nolinkurl{8c9756b28d64dab099da31a4c09229a9e6a2ef35}. Its only import is
\texttt{Std}; it has no Mathlib dependency, no added axioms, no unfinished
proofs, and no use of \texttt{native\_decide}.  The complete compiler output
and a source hash are in \texttt{evidence/lean-build.txt} and
\texttt{evidence/lean-build.json}.

Thirteen named theorems have been checked:
\begin{enumerate}
\item \texttt{anyLE\_correct}: a Boolean scan of $0,\ldots,n$ succeeds
  exactly when a successful index exists in that interval.
\item \texttt{findLE\_sound} and \texttt{findLE\_none\_iff}: the bounded
  witness search returns a valid index, and reports failure exactly when
  the Boolean scan reports failure.
\item \texttt{noBooleanDescent}: an explicitly iterated Boolean condition
  cannot persist under a strict decrease of a natural-number rank.
\item \texttt{quadratic\_expansion}, \texttt{second\_identity}, and
  \texttt{second\_on\_conic}: the quadratic expansion and second-intersection
  identities for an arbitrary integral ternary quadratic form.
\item \texttt{circle\_identity} and \texttt{circle\_inverse\_identity}:
  the polynomial identities underlying the circle parametrization.
\item \texttt{elliptic\_triangle\_identity},
  \texttt{elliptic\_triangle\_area},
  \texttt{triangle\_curve\_certificate}, and
  \texttt{triangle\_on\_curve}: the cleared-denominator identities relating
  a rational right triangle of area $n$ to $y^2=x^3-n^2x$.
\end{enumerate}

The output of \texttt{\#print axioms} lists only
\texttt{propext} and \texttt{Quot.sound}; for
\texttt{elliptic\_triangle\_area} it lists only \texttt{propext}.
In particular, \texttt{Classical.choice} and \texttt{sorryAx} do not occur
in any of the thirteen reports.  The polynomial proofs use a short local
tactic which expands integer powers and products by ring laws, puts
monomials into a fixed order, and asks \texttt{omega} to check the resulting
linear integer identity.  This avoids importing a classical polynomial
normalization theorem merely for the sake of convenient automation.

These are \emph{ordinary Lean theorems about integers and natural numbers}.
They are not yet proof objects of the deeply embedded theory \iPRA.
For example, the source proof of \texttt{anyLE\_correct} uses an ambient
induction whose displayed proposition contains an existential quantifier.
Its object-theory version must instead use the quantifier-free scan
invariants described below.  The absence of classical choice in the
ambient proofs is useful evidence, but does not by itself establish the
induction restriction or produce an \iPRA\ derivation.

\subsection{The existing ProveIt interface}

The inspected ProveIt snapshot is commit
\texttt{b5a3ad037723c44d69a3d5b28a950c554327a989}.
The relevant declarations are in the namespace
\texttt{LeanProofs.ConciseIPRA}.  The file \texttt{PR.lean} defines:
\begin{itemize}
\item \texttt{PR k}, the syntax of $k$-ary primitive recursive function
  symbols, with zero, successor, projections, composition, and primitive
  recursion;
\item \texttt{PTm n}, terms in $n$ free variables;
\item \texttt{Fm PTm n}, first-order intuitionistic formulas;
\item \texttt{PRAAx}, the defining equations, successor axioms, and
  induction for formulas certified quantifier-free;
\item \texttt{PRA.qfInd}, the derived rule that constructs precisely such
  an induction proof.
\end{itemize}
The theorem \texttt{pra\_iff\_native} in
\texttt{NativeBiInterpretation.lean} has the exact form
\[
 (\Gamma\vdash[\mathtt{PRAAx}]\varphi)
 \quad\longleftrightarrow\quad
 \mathtt{NPrf}\bigl(\Gamma.\mathtt{map}\ \mathtt{NFm.ofPRA}\bigr)
                     (\mathtt{NFm.ofPRA}\ \varphi).
\]
Thus the principal completion target for each theorem of this article is
an actual term of
\[
 []\vdash[\mathtt{PRAAx}]\varphi_G,
 \qquad \varphi_G:\mathtt{Fm}\ \mathtt{PTm}\ 0,
\]
where $\varphi_G$ is its explicitly encoded arithmetic statement.
Once that term is constructed, applying the forward direction of
\texttt{pra\_iff\_native} produces the corresponding native-scalar proof.
The bridge is not a rule turning any true theorem of Lean into a theorem
of \iPRA; it takes an \iPRA\ derivation as its input.

\subsection{A concrete module plan}

The following dependency order separates implementation work from
mathematical gaps.  The human proofs of the requisite lemmas are given in
the preceding chapters; the tasks here are their syntactic realization.

\paragraph{1. Encodings and primitive recursive symbols.}
Create \path{DG/Encoding.lean} and \path{DG/PRSymbols.lean}.
Use the signed-integer, rational, vector, polynomial, and finite-list
encodings fixed in the foundations chapter.  Every externally visible
operation must return a total code, with an explicit failure flag where
an input condition fails.  In particular, division by zero, a zero
projective tuple, and a malformed certificate must not produce an
unproved obligation in the implementation.  Define bounded division,
bounded gcd search, exact quotient tests, integer square-root search,
finite products, and coordinate scans as explicit combinations of
\texttt{PR.comp} and \texttt{PR.prec}.

The required interface is not merely a Lean function
\texttt{Nat -> Nat}; it consists of a code in \texttt{PR k}, its defining
equations as terms, and derivations of the equations used by subsequent
proofs.  Each fixed algorithm has such a code.  No universal total
primitive recursive evaluator for arbitrary primitive recursive programs
is assumed.  The uniform algorithms used here evaluate bounded polynomial
data and finite certificates, which is a different and much weaker task.

\paragraph{2. Elementary integer arithmetic inside the theory.}
Create \path{DG/Arithmetic.lean}.  Derive addition, multiplication,
order, and signed arithmetic laws from the defining equations, using
\texttt{PRA.qfInd}.  Prove the bounded division invariants
$a=qb+r$ and $r<b$ for $b>0$, Euclid's divisibility lemma, and the gcd and
coprime square-product lemmas used in descent.  Prime factorization can
be implemented by trial division through the input bound; it needs no
unbounded least-prime operator.  Whenever a proof is phrased with bounded
universal quantification, replace it by the equality-to-$1$ of its
primitive recursive Boolean product before invoking induction.

Create \path{DG/PolynomialCertificates.lean} next.  A reflected
polynomial normalizer should return both a coefficient array and an
object-theory equality proof.  Its soundness is proved by structural
recursion on the polynomial expression and the already derived ring
laws.  This reconstructs the role of the checked Lean identities without
treating an ambient \texttt{Int} identity as an object-theory axiom.
The relatively small identities in \texttt{DiophantineCore.lean} are
regression targets for this compiler.

\paragraph{3. Bounded scans and projective normalization.}
Create \path{DG/FiniteSearch.lean} and \path{DG/Projective.lean}.
For witness search, keep a flag and a candidate index in the recursive
state.  The soundness invariant says that either the flag is false or the
candidate is in the scanned interval and satisfies the Boolean test.
For completeness, fix a proposed successful index $j$ as a numerical
parameter and induct on the scan bound $i$ in the quantifier-free
statement
\[
 j\leq i\ \wedge\ p(j)=1\quad\Longrightarrow\quad
 \operatorname{anyLE}(p,i)=1.
\]
Here $p$ denotes a fixed primitive recursive symbol, not a quantified
function variable of the object language.  Universal generalization over
$j$ is performed after the induction.  Alternatively, maintain the
Boolean product of all failed tests as one additional state component.

Prove canonical projective normalization and uniqueness by gcd
divisibility and a fixed sign convention.  The projective point code
then has literal equality, so searches do not need a quotient type in
the object theory.  Instantiate the scan theorem with the explicit box
size $(2B+1)^{n+1}$ to obtain the constructive Northcott enumeration,
membership correctness, and finite cardinality bound.

\paragraph{4. Heights and polynomial-map certificates.}
Create \path{DG/Height.lean} and \path{DG/Resultant.lean}.
The upper height estimate follows by summing bounded monomial estimates.
For the lower estimate, formalize the two homogeneous B\'ezout
identities supplied by the Sylvester adjugate.  A certificate checker
first compares polynomial coefficients.  A separate soundness theorem
then shows that the common divisor of the two evaluated forms divides
the resultant constant, and combines the certificate's coefficient
bounds with the triangle inequality.

All determinant computations can use finite sums over coded
permutations, with inversion parity computed by bounded loops.
The adjugate identity is a finite polynomial identity, so its proof can
be produced by the same certificate mechanism.  This route does not
require an imported, unbounded algebraic-closure or Nullstellensatz
theorem.  The nonzero determinant test is an integer comparison and the
certificate coefficients are obtained by a finite explicit algorithm.
For a theorem uniform in the matrix size, also formalize cofactor expansion
and the sign-reversing permutation pairing as bounded-sum identities with
quantifier-free induction. Checking a numerical certificate verifies an
instance; it must not be substituted for this uniform arithmetic lemma.

\paragraph{5. Conic charts.}
Create \path{DG/Conics.lean}.  The checked second-intersection identity
settles only incidence.  The remaining formal lemmas must establish that
the produced tuple is nonzero under the nonsingularity and direction
hypotheses, handle the tangent parameter, prove the inverse chart
equations, and establish projective uniqueness and the stated height
bounds.  Work first with cleared integer coordinates and determinants.
Introduce normalized rational/projective codes only after these
polynomial equalities and nonvanishing guards have been proved.
In this way every division in the geometric prose has an explicit
arithmetical justification.

\paragraph{6. Descent and the elliptic example.}
Create \path{DG/FermatDescent.lean} and
\path{DG/CongruentNumber.lean}. Encode the positive integer solution
data as one natural number and define the descent step by the bounded
gcd, quotient, parity, and square-root operations.  Check separately
that it preserves the solution equations and decreases the displayed
rank.  These are quantifier-free implications about fixed primitive
recursive functions.

The generic descent argument is then compiled from the invariant
\[
 \operatorname{Good}(s_k)=1\quad\wedge\quad
 \operatorname{rank}(s_k)+k\leq\operatorname{rank}(s_0),
 \qquad s_{k+1}=\operatorname{Step}(s_k),
\]
at $k=\operatorname{rank}(s_0)+1$.  This is the exact invariant of the
checked \texttt{noBooleanDescent} theorem.  It uses no minimal
counterexample principle and no induction over an unbounded existential
formula.  Finally compile the rational triangle--elliptic-curve
transformations, including the nonzero-denominator and sign conditions,
and deduce the complete rational-point list for $y^2=x^3-x$.

\paragraph{7. Arithmetic dynamics.}
Create \path{DG/Dynamics.lean}.  Starting from the certified two-sided
height estimate, formalize strict height growth outside the explicit
threshold.  Above that threshold, iterate for a specified number of
steps and use quantifier-free arithmetic invariants to rule out a repeat.
Inside the height box, use the complete projective list and a bounded
orbit table.  Equality and membership are Boolean tests on normalized
codes, so the finite repeat test is primitive recursive.  This proves
both soundness and completeness of the algorithm listing all rational
preperiodic points.  It is important to prove the escape lemma before
discarding an orbit that leaves the box; finite search by itself would
not justify that exclusion.

\paragraph{8. Rational roots, local certificates, and denominator geometry.}
Create \path{DG/RationalRoots.lean} and
\path{DG/LocalObstructions.lean}.  The former needs only denominator
clearing, coprimality of powers, and the bounded divisor enumerator.
The latter proves reduction of polynomial evaluation modulo a supplied
modulus and preservation of unimodularity of a primitive integer tuple.
Its finite negative residue certificate is a Boolean conjunction, not a
local-to-global existence principle.

Create \path{DG/WeierstrassDenominators.lean}.  Prove reduced-fraction
uniqueness and the equation $e^2=q^3$ first.  The bounded factorization
module gives the exponent equality $2v=3u$ for each prime, hence the
square/cube denominator form.  Compile the independent square-test
enumerator and prove its list equivalence to the height predicate.
Finally \path{DG/HeightDescent.lean} can express
Theorem~\ref{thm:generation} as a schema for fixed group operations with
supplied object-theory derivations of the group laws, representative
coverage, and reducing-selector properties.  Its two induction invariants
are the truncated height rank and equation~\eqref{eq:descent-expansion}.
This conditional module must retain those hypotheses rather than silently
instantiating them for an arbitrary elliptic curve.

\paragraph{9. Object-theory theorem objects and native transfer.}
Create \path{DG/Statements.lean}, \path{DG/Proofs.lean}, and
\path{DG/NativeTransfer.lean}.  Write each theorem as a closed formula
over the preceding primitive recursive symbols.  Instantiate the
derived arithmetic and search lemmas to obtain
$[]\vdash[\mathtt{PRAAx}]\varphi_G$.  Run an audit recording every use
of \texttt{PRA.qfInd} and the associated proof of
\texttt{Fm.QF}.  Only then apply \texttt{pra\_iff\_native}.
The final integration test should print the axioms of the object-theory
proof terms and ensure that no claimed geometric conclusion has entered
as an extra axiom or as an assumption to a conditional transfer theorem.

\subsection{Completion criteria}

The ordinary Lean arithmetic core is complete and compiled.  The
projective normalization, full conic bijection, Fermat descent step,
resultant-height argument, and preperiodic-point enumeration have full
human proofs in this article but are not fully formalized in Lean here.
The deeply embedded \iPRA\ derivations of the geometric theorems are
also not supplied.  Their completion requires the modules just listed,
with the explicit encodings and induction invariants given here; it does
not require a new mathematical finiteness theorem.

To reproduce the completed part on a machine with the pinned Lean
toolchain, run
\begin{verbatim}
cd lean
lake build
\end{verbatim}
or run \texttt{lean DiophantineCore.lean} directly.  No network fetch of
Mathlib is needed for this core.  Integration with the ProveIt proof
syntax should use the pinned repository snapshot and its own build
configuration, rather than copying the standalone build's dependency
assumptions to the entire repository.

\section{Executable artifacts and reproducibility}
\label{sec:artifacts}

The archive is self-contained for reading the report, rebuilding the PDF,
running the exact arithmetic examples, and compiling the standalone Lean
core. It does not contain a duplicate checkout of ProveIt or a bundled Lean
toolchain. The pinned upstream revision and file blob identifiers are
recorded in \texttt{evidence/source\_manifest.json}.

\begin{center}
\begin{tabular}{p{0.32\textwidth}p{0.59\textwidth}}
\toprule
Path & Content \\
\midrule
\path{diophantine_geometry_ipra.tex}, \texttt{sections/}
& Complete LaTeX source of the human-readable argument \\
\path{diophantine_geometry_ipra.pdf}
& Compiled report \\
\texttt{code/geometry.py}
& Canonical points, height boxes, rational roots, congruence tests,
resultant certificates, preperiodic-point enumeration, and cubic searches \\
\texttt{code/conics.py}
& General pointed-conic parametrization, inverse, and explicit height constants \\
\texttt{code/descent.py}
& Pythagorean extraction, square-factor arithmetic, the descent map,
and signed triangle--cubic transformations \\
\texttt{code/check\_all.py}
& Reproducible exact verification of the executable companions \\
\texttt{lean/}
& Standalone Lean source, toolchain declaration, and Lake project \\
\texttt{evidence/}
& Compiler output, source provenance, exact computation results,
and the verification summary \\
\bottomrule
\end{tabular}
\end{center}

The Python programs use only the standard library, including exact integer
arithmetic and \texttt{fractions.Fraction}. They are reference implementations,
not a proof checker for all the claims in this paper. On malformed inputs
they may raise \texttt{ValueError}; their object-theory versions use an
explicit validity flag and a fixed default value instead. This distinction
does not affect the valid-input statements proved above.

Run the exact verification from the package root with
\begin{verbatim}
python3 code/check_all.py
python3 code/geometry.py
\end{verbatim}
The first command writes \texttt{evidence/validation.json}. It checks 729
ordered quadratic coefficient pairs with entries in $\{-1,0,1\}$, including
480 nonzero-resultant pairs and 11,520 exact height comparisons. It also
checks 25 elliptic point lists against an independent rational-square
enumerator, 4,800 general-conic parameter cases, 20 independent conic
point-list comparisons, 3,931 coprime square extractions, 1,474 Pythagorean
parameter cases, and 5,896 signed rational correspondences. These finite
checks are useful for detecting implementation errors. Completeness of the
mathematical theorems comes from their proofs, not the number of examples.

With Lean 4.32.0 available, the checked core is rebuilt by
\begin{verbatim}
cd lean
lake build
\end{verbatim}
Alternatively, running \texttt{lean DiophantineCore.lean} directly prints
the axiom reports for all thirteen named theorems. There is no package
download or Mathlib dependency. The relevant source hash and actual
compiler output accompany the report.

The PDF is rebuilt from the root by
\begin{verbatim}
latexmk -pdf -interaction=nonstopmode -halt-on-error \
  diophantine_geometry_ipra.tex
\end{verbatim}
The build requires an ordinary LaTeX installation with the packages listed
in the main source. A supplied \texttt{Makefile} provides the targets
\texttt{pdf}, \texttt{check}, and \texttt{lean}.

Independent reviews of the arithmetic foundations, conic height estimates,
finite descent, and resultant dynamics informed the final text. The
remaining formalization boundary is explicit: constructing all closed
\texttt{Prf PRAAx} values described in Section~\ref{sec:lean} is future
implementation work. No unchecked placeholder is presented as a completed
formal theorem.


\clearpage
\begin{thebibliography}{9}
\bibitem{proveit}
Vladimir Reshetnikov and contributors,
\emph{ProveIt: Concise presentations of intuitionistic primitive recursive arithmetic},
revision \texttt{b5a3ad037723c44d69a3d5b28a950c554327a989}, inspected November 15, 2026.
\url{https://github.com/VladimirReshetnikov/ProveIt/tree/b5a3ad037723c44d69a3d5b28a950c554327a989/Logic/PrimitiveRecursiveArithmetic/ConciseIPRA}.
Specific source files and blob identifiers appear in \texttt{evidence/source\_manifest.json}.

\bibitem{milne}
J. S. Milne, \emph{Elliptic Curves}, second edition, World Scientific, 2021.
Author's online version, especially Chapter I, \S2, and Chapter IV.
\url{https://jmilne.org/math/Books/EC2.pdf}.

\bibitem{sutherland2013}
Andrew V. Sutherland, \emph{18.782 Introduction to Arithmetic Geometry,
Lecture 2}, MIT, Fall 2013.
\url{https://ocw.mit.edu/courses/18-782-introduction-to-arithmetic-geometry-fall-2013/}.

\bibitem{sutherland2025}
Andrew V. Sutherland, \emph{18.783 Elliptic Curves, Lecture 1},
MIT, September 4, 2025, especially the discussion of rational points in genus zero.
\url{https://math.mit.edu/classes/18.783/2025/LectureSlides1.pdf}.

\bibitem{silverman}
Joseph H. Silverman, \emph{The Arithmetic of Dynamical Systems},
Graduate Texts in Mathematics 241, Springer, 2007.
Author's book page: \url{https://www.math.brown.edu/johsilve/ADSHome.html}.

\bibitem{silverman2011}
Joseph H. Silverman, \emph{Geometry and Arithmetic of Dominant Rational
Self-Maps of Projective Space}, Shirahama lecture slides, December 2011,
slides on height functoriality and preperiodic points.
\url{https://www.math.brown.edu/johsilve/Presentations/Japan2011.pdf}.

\bibitem{ConradCongruent}
Keith Conrad, \emph{The Congruent Number Problem}, author's expository notes.
\url{https://kconrad.math.uconn.edu/blurbs/ugradnumthy/congnumber.pdf}.

\bibitem{DelahayeMayero}
David Delahaye and Micaela Mayero,
\emph{Diophantus' 20th Problem and Fermat's Last Theorem for $n=4$:
Formalization of Fermat's Proofs in the Coq Proof Assistant}, 2005.
\url{https://arxiv.org/abs/cs/0510011}.
\end{thebibliography}
\end{document}
